% Time-stamp: <2026-01-10 14:49:13 komori>
% Copyright (c) 2021 Yasushi Komori
% All rights reserved.

%%%%%%%%%%%%%%%%%%%%%%%%%%%%%%%%%%%%%%%%%%%%%%%%%%%%%%%%%%%%%%%%
%
% pgfsys
%


% output : 
%
% \brl@pgf@size_x_y : size at (x,y)
% \brl@pgf@minxaty_y : min of x at y
% \brl@pgf@maxxaty_y : max of x at y
% \brl@pgf@minx : min of x
% \brl@pgf@miny : min of y
% \brl@pgf@maxx : max of x
% \brl@pgf@maxy : max of y
%
% context : \empty or matrix_xxx_yyy_
% 
% in draw
%
% \brl@pgf@nodeboxa : box of nodes
% \brl@pgf@nodeboxa : box of nodes
%
% in draw matrix
%
% \brl@pgf_matrix@xxx_yyy a
% \brl@pgf_matrix@xxx_yyy b
%
% in plot
% \brl@pgf@nodetoks : toks of nodes


% \pgf@xb, \pgf@yb : start point
% \pgf@xc, \pgf@yc : current point
%
\newbox\brl@pgf@nodeboxa
\newbox\brl@pgf@nodeboxb

\newtoks\brl@pgf@invoketoks%
\newtoks\brl@pgf@invokemattoks%

\newtoks\brl@pgf@nodetoks%
\def\brl@pgf_appendnodetoks#1{\global\brl@pgf@nodetoks\expandafter{\the\brl@pgf@nodetoks#1}}


\newtoks\brl@pgf@settoks%
\def\brl@pgf_appendsettoks#1{\brl@pgf@settoks\expandafter{\the\brl@pgf@settoks#1}}


% Path construction:
\protected\def\brl@pgfsys@moveto#1#2{\pgfsysprotocol@literal{\brl@pgf_moveto{#1}{#2}}}
\protected\def\brl@pgfsys@lineto#1#2{\pgfsysprotocol@literal{\brl@pgf_lineto{#1}{#2}}}
\protected\def\brl@pgfsys@curveto#1#2#3#4#5#6{\pgfsysprotocol@literal{\brl@pgf_curveto{#1}{#2}{#3}{#4}{#5}{#6}}}
\protected\def\brl@pgfsys@rect#1#2#3#4{\pgfsysprotocol@literal{\brl@pgf_rect{#1}{#2}{#3}{#4}}}
\protected\def\brl@pgfsys@closepath{\pgfsysprotocol@literal{\brl@pgf_closepath}}

% Path usage:
\protected\def\brl@pgfsys@stroke{\pgfsysprotocol@literal{\brl@pgf_stroke}}%
\protected\def\brl@pgfsys@closestroke{\pgfsysprotocol@literal{\brl@pgf_closepath\brl@pgf_stroke}}
\protected\def\brl@pgfsys@fill{\pgfsysprotocol@literal{%
    \brl@pgf_fill%
    {\ifx\tikz@textcolor\empty\tikz@fillcolor\else\tikz@textcolor\fi}%
    {\ifpgfsys@eorule\brl@qs@eorule\else\brl@qs@zerorule\fi}%
  }%
}%
\protected\def\brl@pgfsys@fillstroke{\pgfsysprotocol@literal{%
    \brl@pgf_fillstroke%
    {\ifx\tikz@textcolor\empty\tikz@fillcolor\else\tikz@textcolor\fi}%
    {\ifpgfsys@eorule\brl@qs@eorule\else\brl@qs@zerorule\fi}%
  }%
}%
\protected\def\brl@pgfsys@clipnext{\pgfsysprotocol@literal{\brl@pgf_clipnext{\ifpgfsys@eorule\brl@qs@eorule\else\brl@qs@zerorule\fi}}}%

\def\brl@pgfsys@discardpath{}%


% Transformation:
\protected\def\brl@pgfsys@transformcm#1#2#3#4#5#6{%
  \brl@pgf@x#5\relax%
  \brl@pgf@y#6\relax%
  \pgfsysprotocol@literal{\brl@pgf_transformcm{#1}{#2}{#3}{#4}{\the\brl@pgf@x}{\the\brl@pgf@y}}%
}


% Scopes
\protected\def\brl@pgfsys@beginscope{\pgfsysprotocol@literal{\brl@pgf_beginscope}}%
\protected\def\brl@pgfsys@endscope{\pgfsysprotocol@literal{\brl@pgf_endscope}}%


\protected\def\brl@pgf_beginscope{\begingroup}
\protected\def\brl@pgf_endscope{\endgroup}

\protected\def\brl@pgfsys@setdash#1#2{\pgfsysprotocol@literal{\brl@pgf_setdash{#1}{#2}}}%



% \protected\def\brl@pgfsys@setlinewidth#1{\pgfsysprotocol@literal{\brl@pgf_setlinewidth{#1}}}

%
% \pgfsys@setlinewidth is executed in \tikz@options of ``Step 5' : Setup options'' in \tikz@finish. (tikz.code.tex)
%
\newif\ifbrl@pgfsys@double% override \tikz@mode@double
\let\tikz@mode@doublefalse\brl@pgfsys@doublefalse
\let\tikz@mode@doubletrue\brl@pgfsys@doubletrue

\ifdefined\tikzset
\tikzset{every double/.style={line width=0.5pt}}% forcibly set line width.
\fi

\protected\def\brl@pgfsys@setlinewidth#1{%
  \ifbrl@pgfsys@double%
    \pgfsysprotocol@literal{%
      \brl@pgf_setlinewidth_setdouble{#1}%
    }%
  \else%
    \pgfsysprotocol@literal{%
      \brl@pgf_setlinewidth{#1}%
    }%
  \fi%
}


% \protected\def\brl@pgfsys@setlinewidth#1{%
%   \iftikz@mode@double%
%     \tikz@mode@doublefalse% forcibly reset the switch and double lines.
%     \pgfsysprotocol@literal{%
%       \brl@pgf_setlinewidth_setdouble{#1}%
%     }%
%   \else%
%     \pgfsysprotocol@literal{%
%       \brl@pgf_setlinewidth{#1}%
%     }%
%   \fi%
% }




\def\brl@pgfsys@setmiterlimit#1{}
\def\brl@pgfsys@buttcap{}
\def\brl@pgfsys@roundcap{}
\def\brl@pgfsys@rectcap{}
\def\brl@pgfsys@miterjoin{}
\def\brl@pgfsys@roundjoin{}
\def\brl@pgfsys@beveljoin{}

% \def\brl@pgfsys@color@rgb@stroke#1#2#3{\pgfsysprotocol@literal{\brl@pgf_stroke}}
% \def\brl@pgfsys@color@cmyk@stroke#1#2#3#4{\pgfsysprotocol@literal{\brl@pgf_stroke}}
% \def\brl@pgfsys@color@cmy@stroke#1#2#3{\pgfsysprotocol@literal{\brl@pgf_stroke}}
% \def\brl@pgfsys@color@gray@stroke#1{\pgfsysprotocol@literal{\brl@pgf_stroke}}

% \def\brl@pgfsys@color@rgb@fill#1#2#3{\pgfsysprotocol@literal{\brl@pgf_color_rgb_fill{#1}{#2}{#3}}}
% \def\brl@pgfsys@color@cmyk@fill#1#2#3#4{\pgfsysprotocol@literal{\brl@pgf_color_cmyk_fill{#1}{#2}{#3}{#4}}}
% \def\brl@pgfsys@color@cmy@fill#1#2#3{\pgfsysprotocol@literal{\brl@pgf_color_cmy_fill{#1}{#2}{#3}}}
% \def\brl@pgfsys@color@gray@fill#1{\pgfsysprotocol@literal{\brl@pgf_color_gray_fill{#1}}}

\def\brl@pgfsys@color@rgb@stroke#1#2#3{}
\def\brl@pgfsys@color@cmyk@stroke#1#2#3#4{}
\def\brl@pgfsys@color@cmy@stroke#1#2#3{}
\def\brl@pgfsys@color@gray@stroke#1{}
\def\brl@pgfsys@color@rgb@fill#1#2#3{}
\def\brl@pgfsys@color@cmyk@fill#1#2#3#4{}
\def\brl@pgfsys@color@cmy@fill#1#2#3{}
\def\brl@pgfsys@color@gray@fill#1{}


\def\brl@pgfsys@setmiterlimit#1{}


% Visibility
\def\brl@pgfsys@begininvisiblescope{}
\def\brl@pgfsys@endinvisiblescope{}

% Opacity
\def\brl@pgfsys@stroke@opacity#1{}%
\def\brl@pgfsys@fill@opacity#1{}%

% Blending
\def\brl@pgfsys@blend@mode#1{}%

% Literals:
\def\brl@pgfsys@invoke#1{%
  \edef\brl@pgfsysprotocol@tmp{{\the\brl@pgf@invoketoks#1}}%
  \global\brl@pgf@invoketoks\brl@pgfsysprotocol@tmp%
}%

\def\brl@pgfsys@invoke_matrix#1{%
  \edef\brl@pgfsysprotocol@tmp{{\the\brl@pgf@invokemattoks#1}}%
  \global\brl@pgf@invokemattoks\brl@pgfsysprotocol@tmp%
}%

\def\brl@pgfsysprotocol@literalbuffered#1{%
  \edef\pgfsysprotocol@temp{{#1}}% <- here \space is removed.
  \expandafter\pgfutil@g@addto@macro\expandafter\pgfsysprotocol@currentprotocol\pgfsysprotocol@temp%
}

% Text:
% output:
% \brl@pgf@nodetoks{\brl@pgf_hbox{xxx}{yyy}{brailles}...}
%
\def\brl@pgfsys@hbox#1{% #1 box : boxname is used
  \pgfsys@begin@idscope%
  % 
    %
  \ifcsname brl@pgf@nodetext_\string#1\endcsname%
    %
    %
    \pgfsysprotocol@literal{\brl@pgf_put{\the\brl@pgf@pt_x}{\the\brl@pgf@pt_y}{\csname brl@pgf@nodetext_\string#1\endcsname}}%
    \global\expandafter\let\csname brl@pgf@nodetext_\string#1\endcsname\@undefined% clear
    \fi%
    %
    % 
    % boxname is of the form \pgfnodepart[boxname]box
    % 
    %% box#1 is ignored. all data is restored in \brl@pgf@nodetext_boxname, which is globally defined in brl_print.
    % 
  \pgfsys@end@idscope%
}

%
% the following is the original code of \pgfsys@hboxsynced, which
% is overwritten by pgfsys-dvipdfmx.
%

\def\brl@pgfsys@hboxsynced#1{%
  \pgfsys@beginscope%
  % \pgflowlevelsynccm
  \pgfsys@transformcm%
  % \pgf_transformcm%
  {\pgf@pt@aa}{\pgf@pt@ab}%
  {\pgf@pt@ba}{\pgf@pt@bb}%
  {\pgf@pt@x}{\pgf@pt@y}%
  \pgfsys@hbox#1%
  \pgfsys@endscope%
}%


%% replace the codes by those for plots.
\def\brl@pgfsys_init{%
  \let\pgfsys@moveto\brl@pgfsys@moveto%
  \let\pgfsys@lineto\brl@pgfsys@lineto%
  \let\pgfsys@curveto\brl@pgfsys@curveto%
  \let\pgfsys@rect\brl@pgfsys@rect%
  \let\pgfsys@closepath\brl@pgfsys@closepath%
  \let\pgfsys@stroke\brl@pgfsys@stroke%
  \let\pgfsys@closestroke\brl@pgfsys@closestroke%
  \let\pgfsys@fill\brl@pgfsys@fill%
  \let\pgfsys@fillstroke\brl@pgfsys@fillstroke%
  \let\pgfsys@clipnext\brl@pgfsys@clipnext%
  \let\pgfsys@discardpath\brl@pgfsys@discardpath%
  \let\pgfsys@transformcm\brl@pgfsys@transformcm%
  \let\pgfsys@beginscope\brl@pgfsys@beginscope%
  \let\pgfsys@endscope\brl@pgfsys@endscope%
  \let\pgfsys@setdash\brl@pgfsys@setdash%
  \let\pgfsys@setmiterlimit\brl@pgfsys@setmiterlimit%
  \let\pgfsys@buttcap\brl@pgfsys@buttcap%
  \let\pgfsys@roundcap\brl@pgfsys@roundcap%
  \let\pgfsys@rectcap\brl@pgfsys@rectcap%
  \let\pgfsys@miterjoin\brl@pgfsys@miterjoin%
  \let\pgfsys@roundjoin\brl@pgfsys@roundjoin%
  \let\pgfsys@beveljoin\brl@pgfsys@beveljoin%
  \let\pgfsys@color@rgb@stroke\brl@pgfsys@color@rgb@stroke%
  \let\pgfsys@color@rgb@fill\brl@pgfsys@color@rgb@fill%
  \let\pgfsys@color@cmyk@stroke\brl@pgfsys@color@cmyk@stroke%
  \let\pgfsys@color@cmyk@fill\brl@pgfsys@color@cmyk@fill%
  \let\pgfsys@color@cmy@stroke\brl@pgfsys@color@cmy@stroke%
  \let\pgfsys@color@cmy@fill\brl@pgfsys@color@cmy@fill%
  \let\pgfsys@color@gray@stroke\brl@pgfsys@color@gray@stroke%
  \let\pgfsys@color@gray@fill\brl@pgfsys@color@gray@fill%
  \let\pgfsys@setmiterlimit\brl@pgfsys@setmiterlimit%
  \let\pgfsys@setlinewidth\brl@pgfsys@setlinewidth%
  \let\pgfsys@begininvisiblescope\brl@pgfsys@begininvisiblescope%
  \let\pgfsys@endinvisiblescope\brl@pgfsys@endinvisiblescope%
  \let\pgfsys@stroke@opacity\brl@pgfsys@stroke@opacity%
  \let\pgfsys@fill@opacity\brl@pgfsys@fill@opacity%
  \let\pgfsys@blend@mode\brl@pgfsys@blend@mode%
  \let\pgfsys@invoke\brl@pgfsys@invoke%
  \let\pgfsys@hbox\brl@pgfsys@hbox%
  \let\pgfsys@hboxsynced\brl@pgfsys@hboxsynced%
  %
  % space after commands should be ignored 
  \let\pgfsysprotocol@literalbuffered\brl@pgfsysprotocol@literalbuffered%
}



%%%%%%%%%%%%%%%%%%%%%%%%%%%%%%%%%%%%%%%%%%%%%%%%%%%%%%%%%%%%%%%%
%
% pgf subroutine
%

\newtoks\brl@pgf@pathtoks
\def\brl@pgf_appendpathtoks#1{\brl@pgf@pathtoks\expandafter{\the\brl@pgf@pathtoks#1}}%

\newtoks\brl@pgf@pointtoks
\def\brl@pgf_appendpointtoks#1{\brl@pgf@pointtoks\expandafter{\the\brl@pgf@pointtoks#1}}%

\newtoks\brl@pgf@cliptoks
\protected\def\brl@pgf_appendcliptoks#1{\brl@pgf@cliptoks\expandafter{\the\brl@pgf@cliptoks#1}}%

\newtoks\brl@pgf@clipytoks
\def\brl@pgf_appendclipytoks#1{\brl@pgf@clipytoks\expandafter{\the\brl@pgf@clipytoks#1}}%

\newtoks\brl@pgf@clipalltoks
\def\brl@pgf_appendclipalltoks#1{\brl@pgf@clipalltoks\expandafter{\the\brl@pgf@clipalltoks#1}}%

\newcount\brl@pgf@lineerr
\newcount\brl@pgf@linedx
\newcount\brl@pgf@linedy
\newcount\brl@pgf@lineex
\newcount\brl@pgf@lineey
\newcount\brl@pgf@dx
\newcount\brl@pgf@dy
\newcount\brl@pgf@ptx
\newcount\brl@pgf@pty
\newcount\brl@pgf@dnum

\newcount\brl@pgf@currx
\newcount\brl@pgf@curry
\newcount\brl@pgf@currxa
\newcount\brl@pgf@currya
\newcount\brl@pgf@currxb
\newcount\brl@pgf@curryb

\newcount\brl@pgf@prevx
\newcount\brl@pgf@prevy
\newcount\brl@pgf@prevxa
\newcount\brl@pgf@prevya
\newcount\brl@pgf@prevxb
\newcount\brl@pgf@prevyb

\newcount\brl@pgf_stroke@pattern_num
\newcount\brl@pgf_stroke@pattern_phase

\newcount\brl@pgf_fill@x
\newcount\brl@pgf_fill@y
\newcount\brl@pgf_fill@num
\newcount\brl@pgf_fill@mody

\newcount\brl@pgf_clip_save@y

\newcount\brl@pgf_filldef@currx%
\newcount\brl@pgf_filldef@curry%
\newcount\brl@pgf_filldef@prevx%
\newcount\brl@pgf_filldef@dotnum%


% local variables
\newdimen\brl@pgf@x
\newdimen\brl@pgf@y
\newdimen\brl@pgf@xa
\newdimen\brl@pgf@ya
\newdimen\brl@pgf@xb
\newdimen\brl@pgf@yb
\newdimen\brl@pgf@xc
\newdimen\brl@pgf@yc
\newdimen\brl@pgf@xd
\newdimen\brl@pgf@yd
\newdimen\brl@pgf@xe
\newdimen\brl@pgf@ye
\newdimen\brl@pgf@xf
\newdimen\brl@pgf@yf
\newdimen\brl@pgf@length

\newdimen\brl@pgf@pt_x
\newdimen\brl@pgf@pt_y

%%%%%%%%%%%%%%%%%%%%%%%%%%%%%%%%%%%%%%%%%%%%%%%%%%%%%%%%%%%%%%%%
%
% path construction
%


\protected\def\brl@pgf_moveto#1#2{%
  \brl@pgf_appendpathtoks{\brl@pgf_moveto_do{#1}{#2}}%
}

\def\brl@pgf_moveto_do#1#2{%
  % These two store the current position:
  \pgf@xc#1\relax%
  \pgf@yc#2\relax%
  % These two store the last moveto:
  \pgf@xb\pgf@xc%
  \pgf@yb\pgf@yc%
  % 
  \brl@pgf@prevx\@M% initialize
  \brl@pgf@prevy\@M%
  \brl@pgf@prevxa\@M% initialize
  \brl@pgf@prevya\@M%
  \brl@pgf@prevxb\@M% initialize
  \brl@pgf@prevyb\@M%
  %
}

\protected\def\brl@pgf_lineto#1#2{%
  \brl@pgf_appendpathtoks{\brl@pgf_lineto_do{#1}{#2}}%
}



\protected\def\brl@pgf_lineto_do#1#2{%
  %
  \brl@pgf@xa#1\relax%
  \brl@pgf@ya#2\relax%
  % 
  \brl@pgf_transform{\brl@pgf@xa}{\brl@pgf@ya}%
  \brl@pgf_transform{\pgf@xc}{\pgf@yc}%
  % 
  \edef\brl@pgf@endx{\the\numexpr\brl@pgf@xa/\p@\relax}% end
  \edef\brl@pgf@endy{\the\numexpr\brl@pgf@ya/\p@\relax}% end
  % 
  \brl@pgf@currx\numexpr\pgf@xc/\p@\relax% start
  \brl@pgf@curry\numexpr\pgf@yc/\p@\relax% start
  % 
  \brl@pgf@linedx\brl@pgf@endx\relax\advance\brl@pgf@linedx-\brl@pgf@currx%
  \brl@pgf@linedy\brl@pgf@endy\relax\advance\brl@pgf@linedy-\brl@pgf@curry%
  %
  % \ifnum\brl@pgf@linedx=\z@%
  %   \ifnum\brl@pgf@linedy=\z@%
  \ifcase\brl@pgf@linedx%
    \ifcase\brl@pgf@linedy%
    \else%
      \brl@pgf_lineto_do_sub%
    \fi%
  \else%
    \brl@pgf_lineto_do_sub%
  \fi%
  \pgf@xc#1\relax%
  \pgf@yc#2\relax%
}


\def\brl@pgf_lineto_do_normalize_double{%
  \pgf@x\brl@pgf@xa\advance\pgf@x-\pgf@xc%
  \pgf@y\brl@pgf@ya\advance\pgf@y-\pgf@yc%
  \brl@pgf_normalize%
}

\def\brl@pgf_lineto_do_sub{%
  % linewidth
  \brl@pgf_lineto_do_normalize%
  %
  % 
  % 
  \brl@pgf_renew_localminmax% currx, curry
  % 
  % 
  \brl@pgf@dx\brl@pgf@prevx\advance\brl@pgf@dx-\brl@pgf@currx%
  \brl@pgf@dy\brl@pgf@prevy\advance\brl@pgf@dy-\brl@pgf@curry%
  \brl@pgf_clip_save_do%
  \brl@pgf@prevx\brl@pgf@currx% start point
  \brl@pgf@prevy\brl@pgf@curry% start point
  %
  \ifx\brl@pgf_stroke_save_do\empty%
  \else%
    \ifnum\brl@pgf@linedx>\z@%
      \let\brl@pgf@onex\@ne%
    \else%
      \brl@pgf@linedx-\brl@pgf@linedx%
      \let\brl@pgf@onex\m@ne%
    \fi%
    %
    \ifnum\brl@pgf@linedy>\z@%
      \brl@pgf@linedy-\brl@pgf@linedy%
      \let\brl@pgf@oney\@ne%
    \else%
      \let\brl@pgf@oney\m@ne%
    \fi%
    %
    \brl@pgf@lineex\brl@pgf@linedx\multiply\brl@pgf@lineex\tw@% ex=2dx
    \brl@pgf@lineey\brl@pgf@linedy\multiply\brl@pgf@lineey\tw@% ey=2dy
    \brl@pgf@lineerr\brl@pgf@lineex\advance\brl@pgf@lineerr\brl@pgf@lineey% ex+ey
    \def\brl@pgf_lineto_next{\brl@pgf_lineto_loop}%
    % 
    %
    %
    %
    \brl@pgf_lineto_loop%
    %
  \fi%
  %
  %
  \brl@pgf@currx\brl@pgf@endx\relax%
  \brl@pgf@curry\brl@pgf@endy\relax%
  % 
  \brl@pgf@dx\brl@pgf@prevx\advance\brl@pgf@dx-\brl@pgf@currx%
  \brl@pgf@dy\brl@pgf@prevy\advance\brl@pgf@dy-\brl@pgf@curry%
  \brl@pgf_clip_save_do%
  \brl@pgf@prevx\brl@pgf@currx%
  \brl@pgf@prevy\brl@pgf@curry%
  %
  % 
  %
}

% only for stroke
\def\brl@pgf_lineto_loop{%
  % 
  \brl@pgf_setpoint_line%
  %
  %
  %
  \ifnum\brl@pgf@endx=\brl@pgf@currx%
    \ifnum\brl@pgf@endy=\brl@pgf@curry%
      \let\brl@pgf_lineto_next\empty%
    \fi%
  \fi%
  \count@\brl@pgf@lineerr%
  \ifnum\count@>\brl@pgf@linedy%
    \advance\brl@pgf@lineerr\brl@pgf@lineey%
    \advance\brl@pgf@currx\brl@pgf@onex%
  \fi%
  \ifnum\count@<\brl@pgf@linedx%
    \advance\brl@pgf@lineerr\brl@pgf@lineex%
    \advance\brl@pgf@curry\brl@pgf@oney%
  \fi%
  \brl@pgf_lineto_next%
}


% for line
\def\brl@pgf_setpoint_line_single{%
  %
  % 
  \brl@pgf@dx\brl@pgf@prevxa\advance\brl@pgf@dx-\brl@pgf@currx%
  \brl@pgf@dy\brl@pgf@prevya\advance\brl@pgf@dy-\brl@pgf@curry%
  % 
  \ifnum\numexpr\brl@pgf@dx*\brl@pgf@dx+\brl@pgf@dy*\brl@pgf@dy\relax>\brl@pgf@leastdist%
    % 
    % 
    % \brl@pgf@prevx, \brl@pgf@currx
    % \brl@pgf@prevy, \brl@pgf@curry
    % \brl@pgf@dx, \brl@pgf@dy
    %
    \brl@pgf_stroke_save_do%
    \brl@pgf@prevxa\brl@pgf@currx%
    \brl@pgf@prevya\brl@pgf@curry%
    % 
  \fi%
}

\def\brl@pgf_setpoint_line_double{%
  %
  % 
  \brl@pgf@currxa\brl@pgf@currx\advance\brl@pgf@currxa\brl@pgf@normal_x\relax%
  \brl@pgf@currya\brl@pgf@curry\advance\brl@pgf@currya\brl@pgf@normal_y\relax%
  % 
  \brl@pgf@dx\brl@pgf@prevxa\advance\brl@pgf@dx-\brl@pgf@currxa%
  \brl@pgf@dy\brl@pgf@prevya\advance\brl@pgf@dy-\brl@pgf@currya%
  % 
  \ifnum\numexpr\brl@pgf@dx*\brl@pgf@dx+\brl@pgf@dy*\brl@pgf@dy\relax>\brl@pgf@leastdist%
    % 
    % 
    % \brl@pgf@prevx, \brl@pgf@currx
    % \brl@pgf@prevy, \brl@pgf@curry
    % \brl@pgf@dx, \brl@pgf@dy
    %
    \brl@pgf@currxb\brl@pgf@currx\advance\brl@pgf@currxb-\brl@pgf@normal_x\relax%
    \brl@pgf@curryb\brl@pgf@curry\advance\brl@pgf@curryb-\brl@pgf@normal_y\relax%
    % 
    \brl@pgf_stroke_save_do%
    \brl@pgf@prevxa\brl@pgf@currxa%
    \brl@pgf@prevya\brl@pgf@currya%
    \brl@pgf@prevxb\brl@pgf@currxb%
    \brl@pgf@prevyb\brl@pgf@curryb%
    %
    % 
  \fi%
}



\protected\def\brl@pgf_curveto#1#2#3#4#5#6{%
  \brl@pgf_appendpathtoks{\brl@pgf_curveto_do{#1}{#2}{#3}{#4}{#5}{#6}}%
}

\def\brl@pgf_curveto_do#1#2#3#4#5#6{%
  % 
  \brl@pgf@xa#1\relax%
  \brl@pgf@ya#2\relax%
  \brl@pgf@xb#3\relax%
  \brl@pgf@yb#4\relax%
  \brl@pgf@xc#5\relax%
  \brl@pgf@yc#6\relax%
  % 
  \brl@pgf_transform{\pgf@xc}{\pgf@yc}%
  \brl@pgf_transform{\brl@pgf@xa}{\brl@pgf@ya}%
  \brl@pgf_transform{\brl@pgf@xb}{\brl@pgf@yb}%
  \brl@pgf_transform{\brl@pgf@xc}{\brl@pgf@yc}%
  %
  % 
  % 
  \pgfutil@tempdima\z@%
  %
  \def\brl@pgf@xasign{0}%
  \brl@pgf_curveto_loop%
  % 
  \pgf@xc#5\relax%
  \pgf@yc#6\relax%
}


\def\brl@pgf_curveto_do_normalize_double{%
  \brl@pgf_normalize%
  \@tempdimb\expandafter\brl@@strip@pt\the\@tempdimc\@tempdimb%
}%


\def\brl@pgf_curveto_loop{%
  % 
  \pgfutil@tempdimb\p@\advance\pgfutil@tempdimb-\pgfutil@tempdima% 1-t
  % 
  \edef\brl@pgf@time_t{\pgf@sys@tonumber{\pgfutil@tempdima}}%
  \edef\brl@pgf@time_s{\pgf@sys@tonumber{\pgfutil@tempdimb}}%
  % 
  %
  % First iteration:
  \brl@pgf@xd\brl@pgf@time_s\pgf@xc\advance\brl@pgf@xd\brl@pgf@time_t\brl@pgf@xa% (1-t) P0 + t P1
  \brl@pgf@yd\brl@pgf@time_s\pgf@yc\advance\brl@pgf@yd\brl@pgf@time_t\brl@pgf@ya%
  % 
  \brl@pgf@xe\brl@pgf@time_s\brl@pgf@xa\advance\brl@pgf@xe\brl@pgf@time_t\brl@pgf@xb% (1-t) P1 + t P2
  \brl@pgf@ye\brl@pgf@time_s\brl@pgf@ya\advance\brl@pgf@ye\brl@pgf@time_t\brl@pgf@yb%
  % 
  \brl@pgf@xf\brl@pgf@time_s\brl@pgf@xb\advance\brl@pgf@xf\brl@pgf@time_t\brl@pgf@xc% (1-t) P2 + t P3
  \brl@pgf@yf\brl@pgf@time_s\brl@pgf@yb\advance\brl@pgf@yf\brl@pgf@time_t\brl@pgf@yc%
  % Second iteration:
  \brl@pgf@xd\brl@pgf@time_s\brl@pgf@xd\advance\brl@pgf@xd\brl@pgf@time_t\brl@pgf@xe% (1-t)^2 P0 + 2(1-t)t P1 + t^2 P2
  \brl@pgf@yd\brl@pgf@time_s\brl@pgf@yd\advance\brl@pgf@yd\brl@pgf@time_t\brl@pgf@ye%
  % 
  \brl@pgf@xe\brl@pgf@time_s\brl@pgf@xe\advance\brl@pgf@xe\brl@pgf@time_t\brl@pgf@xf% (1-t)^2 P1 + 2(1-t)t P2 + t^2 P3
  \brl@pgf@ye\brl@pgf@time_s\brl@pgf@ye\advance\brl@pgf@ye\brl@pgf@time_t\brl@pgf@yf%
  % Third iteration:
  \brl@pgf@xf\brl@pgf@time_s\brl@pgf@xd\advance\brl@pgf@xf\brl@pgf@time_t\brl@pgf@xe%
  \brl@pgf@yf\brl@pgf@time_s\brl@pgf@yd\advance\brl@pgf@yf\brl@pgf@time_t\brl@pgf@ye%
  %
  \pgf@x\brl@pgf@xe\advance\pgf@x-\brl@pgf@xd% derivative / 3
  \pgf@y\brl@pgf@ye\advance\pgf@y-\brl@pgf@yd% derivative / 3
  %
  \brl@pgf_curveto_do_normalize%
  % 
  \brl@pgf@currx\numexpr\brl@pgf@xf/\p@\relax%
  \brl@pgf@curry\numexpr\brl@pgf@yf/\p@\relax%
  %
  %
  \brl@pgf_setpoint_curve%
  %
  %
  % determine the next step
  %
  % f(t)  = (1-t)^3 P0 + 3(1-t)^2t P1 + 3(1-t)t^2 P2 + t^3 P3
  % f'(t) = 3( (P1-P0)(1-t)^2 + 2(P2-P1)(1-t)t + (P3-P2)t^2 )
  %
  \ifx\brl@pgf@xasign\brl@pgf@xsign%
    \ifx\brl@pgf@yasign\brl@pgf@ysign%
    \else% min or max
      \brl@pgf_renew_localminmax%
    \fi%
  \else% min or max
    \brl@pgf_renew_localminmax%
  \fi%
  \let\brl@pgf@xasign\brl@pgf@xsign%
  \let\brl@pgf@yasign\brl@pgf@ysign%
  %
  \ifdim\@tempdimb>512\p@%
    \brl@pgf@length0.000651042\p@% 1/512/3 = 0.00195313/3
  \else%
    \ifdim\@tempdimb>3.33333\p@%
      \brl@pgf@length\dimexpr0.333333\p@*\p@/\@tempdimb\relax%
    \else%
      \brl@pgf@length0.1\p@% 
    \fi%
  \fi%
  %
  % 
  \advance\pgfutil@tempdima\brl@pgf@length% almost 1
  \ifdim\p@<\pgfutil@tempdima%
  \else%
    \expandafter\brl@pgf_curveto_loop%
  \fi%
}




\def\brl@pgf_setpoint_curve_single{%
  %
  %
  \brl@pgf@dx\brl@pgf@prevx\advance\brl@pgf@dx-\brl@pgf@currx%
  \brl@pgf@dy\brl@pgf@prevy\advance\brl@pgf@dy-\brl@pgf@curry%
  % 
  \ifnum\numexpr\brl@pgf@dx*\brl@pgf@dx+\brl@pgf@dy*\brl@pgf@dy\relax>\brl@pgf@leastdist%
    % 
    % 
    % \brl@pgf@prevx, \brl@pgf@currx
    % \brl@pgf@prevy, \brl@pgf@curry
    % \brl@pgf@dx, \brl@pgf@dy
    % 
    \brl@pgf_clip_save_do%
    \brl@pgf@prevx\brl@pgf@currx%
    \brl@pgf@prevy\brl@pgf@curry%
    % 
    \brl@pgf_stroke_save_do%
    \brl@pgf@prevxa\brl@pgf@currx%
    \brl@pgf@prevya\brl@pgf@curry%
    % 
  \fi%
}


\def\brl@pgf_setpoint_curve_double{%
  %
  %
  %
  \ifx\brl@pgf_stroke_save_do\empty%
  \else%
    \brl@pgf@currxa\brl@pgf@currx\advance\brl@pgf@currxa\brl@pgf@normal_x\relax%
    \brl@pgf@currya\brl@pgf@curry\advance\brl@pgf@currya\brl@pgf@normal_y\relax%
    \brl@pgf@dx\brl@pgf@prevxa\advance\brl@pgf@dx-\brl@pgf@currxa%
    \brl@pgf@dy\brl@pgf@prevya\advance\brl@pgf@dy-\brl@pgf@currya%
    % 
    \ifnum\numexpr\brl@pgf@dx*\brl@pgf@dx+\brl@pgf@dy*\brl@pgf@dy\relax>\brl@pgf@leastdist%
      % 
      \brl@pgf@currxb\brl@pgf@currx\advance\brl@pgf@currxb-\brl@pgf@normal_x\relax%
      \brl@pgf@curryb\brl@pgf@curry\advance\brl@pgf@curryb-\brl@pgf@normal_y\relax%
      \brl@pgf@dx\brl@pgf@prevxb\advance\brl@pgf@dx-\brl@pgf@currxb%
      \brl@pgf@dy\brl@pgf@prevyb\advance\brl@pgf@dy-\brl@pgf@curryb%
      % 
      \ifnum\numexpr\brl@pgf@dx*\brl@pgf@dx+\brl@pgf@dy*\brl@pgf@dy\relax>\brl@pgf@leastdist%
        \brl@pgf_stroke_save_do%
        \brl@pgf@prevxa\brl@pgf@currxa%
        \brl@pgf@prevya\brl@pgf@currya%
        \brl@pgf@prevxb\brl@pgf@currxb%
        \brl@pgf@prevyb\brl@pgf@curryb%
      \fi%
    \fi%
  \fi%
  %
  \brl@pgf@dx\brl@pgf@prevx\advance\brl@pgf@dx-\brl@pgf@currx%
  \brl@pgf@dy\brl@pgf@prevy\advance\brl@pgf@dy-\brl@pgf@curry%
  \ifnum\numexpr\brl@pgf@dx*\brl@pgf@dx+\brl@pgf@dy*\brl@pgf@dy\relax>\brl@pgf@leastdist%
    % 
    % 
    % \brl@pgf@prevx, \brl@pgf@currx
    % \brl@pgf@prevy, \brl@pgf@curry
    % \brl@pgf@dx, \brl@pgf@dy
    % 
    \brl@pgf_clip_save_do%
    \brl@pgf@prevx\brl@pgf@currx%
    \brl@pgf@prevy\brl@pgf@curry%
    % 
  \fi%
}



\protected\def\brl@pgf_rect#1#2#3#4{%
  \brl@pgf_appendpathtoks{\brl@pgf_rect_do{#1}{#2}{#3}{#4}}%
}

\def\brl@pgf_rect_do#1#2#3#4{%
  \pgf@x#1\relax%
  \pgf@y#2\relax%
  \pgf@xa#3\relax\advance\pgf@xa\pgf@x%
  \pgf@ya#4\relax\advance\pgf@ya\pgf@y%
  \brl@pgf_moveto_do{\pgf@x}{\pgf@y}%
  \brl@pgf_lineto_do{\pgf@xa}{\pgf@y}%
  \brl@pgf_lineto_do{\pgf@xa}{\pgf@ya}%
  \brl@pgf_lineto_do{\pgf@x}{\pgf@ya}%
  \brl@pgf_closepath_do%
}

\protected\def\brl@pgf_closepath{%
  \brl@pgf_appendpathtoks{\brl@pgf_closepath_do}%
}

\def\brl@pgf_closepath_do{%
  \brl@pgf_lineto_do{\pgf@xb}{\pgf@yb}%
  \brl@pgf_clip_save_closepath%
  \brl@pgf_modify% if the closed path is too small, then it is replaced by a big dot.
}


\def\brl@pgf_renew_localminmax{%
  \ifnum\brl@pgf@localminx>\brl@pgf@currx%
    \edef\brl@pgf@localminx{\the\brl@pgf@currx}%
  \fi%
  \ifnum\brl@pgf@localminy>\brl@pgf@curry%
    \edef\brl@pgf@localminy{\the\brl@pgf@curry}%
  \fi%
  \ifnum\brl@pgf@localmaxx<\brl@pgf@currx%
    \edef\brl@pgf@localmaxx{\the\brl@pgf@currx}%
  \fi%
  \ifnum\brl@pgf@localmaxy<\brl@pgf@curry%
    \edef\brl@pgf@localmaxy{\the\brl@pgf@curry}%
  \fi%
}

\def\brl@pgf_modify{%
  \ifnum\brl@pgf@localminx=\@M%
  \else%
    \@tempcnta\brl@pgf@localmaxx\advance\@tempcnta-\brl@pgf@localminx\relax%
    \@tempcntb\brl@pgf@localmaxy\advance\@tempcntb-\brl@pgf@localminy\relax%
    \ifnum\numexpr\@tempcnta*\@tempcnta+\@tempcntb*\@tempcntb\relax<\csname brl@pgf_modify@leastdist\brl@pgf@pointsize\endcsname\relax% very small
      \brl@pgf@pointtoks{}% no storke
      \def\brl@pgf_clip@maxy{-\@M}% no fill or no clip
      \def\brl@pgf_clip@miny{\@M}%
      % set a big point.
      \brl@pgf@currx\numexpr(\brl@pgf@localmaxx+\brl@pgf@localminx)/\tw@\relax%
      \brl@pgf@curry\numexpr(\brl@pgf@localmaxy+\brl@pgf@localminy)/\tw@\relax%
      \brl@pgf@prevx\@M% initialize
      \brl@pgf@prevy\@M%
      \brl@pgf@prevxa\@M% initialize
      \brl@pgf@prevya\@M%
      \brl@pgf@prevxb\@M% initialize
      \brl@pgf@prevyb\@M%
      {%
        \def\brl@pgf@pointsize{3}% a big point
        \brl@pgf_drawpoint%
      }%
    \fi%
  \fi%
}


%%
%%%%%%%%%%%%%%%%%%%%%%%%%%%%%%%%%%%%%%%%%%%%%%%%%%%%%%%%%%%%%%%%
%
% path usage
%


% Graphics state
%
% #1 : pattern. on off on off ...
% #2 : phase of the start.
%
% \brl@pgf_stroke@pattern_name_1pt,2pt,... = \brl@pgf_stroke@pattern_\the\brl@pgf_stroke@pattern_num _0
%
% \brl@pgf_stroke@pattern_\the\brl@pgf_stroke@pattern_num _X
% X = 0 : X = 1 & \brl@pgf_appendpoint
% X = 1 : X = 2
% X = 2 : X = 3 & \brl@pgf_appendpoint etc.
% ...
% X = n equiv X = 0
%
% no pattern
% \brl@pgf_stroke@pattern_name_ = \brl@pgf_appendpoint




\protected\def\brl@pgf_stroke{%
  % stroke_on
  \brl@pgf@pointtoks{}%
  \let\brl@pgf_stroke_save_do\brl@pgf_stroke_save%
  % clip_off
  \let\brl@pgf_clip_save_do\empty%
  %
  %
  \brl@pgf_reset_localminmax%
  \the\brl@pgf@pathtoks%
  \brl@pgf@pathtoks{}%
  %
  % stroke_do
  \brl@pgf_stroke_do%
}


\protected\def\brl@pgf_fill#1#2{%
  \brl@pgf_color{#1}%
  % stroke_off
  \let\brl@pgf_stroke_save_do\empty%
  % clip_on
  \brl@pgf@cliptoks{}%
  \let\brl@pgf_clip_save_do\brl@pgf_clip_save%
  \def\brl@pgf_clip@maxy{-\@M}%
  \def\brl@pgf_clip@miny{\@M}%
  %
  \brl@pgf_reset_localminmax%
  #2% eorule
  \the\brl@pgf@pathtoks%
  \brl@pgf@pathtoks{}%
  \brl@pgf_clip_save_closepath%
  %
  % fill_do
  \brl@pgf_fill_do%
}

\protected\def\brl@pgf_fillstroke#1#2{%
  \brl@pgf_color{#1}%
  % stroke_on 
  \brl@pgf@pointtoks{}%
  \let\brl@pgf_stroke_save_do\brl@pgf_stroke_save%
  % clip_on
  \brl@pgf@cliptoks{}%
  \let\brl@pgf_clip_save_do\brl@pgf_clip_save%
  \def\brl@pgf_clip@maxy{-\@M}%
  \def\brl@pgf_clip@miny{\@M}%
  %
  \brl@pgf_reset_localminmax%
  #2% eorule
  \the\brl@pgf@pathtoks%
  \brl@pgf@pathtoks{}%
  \brl@pgf_clip_save_closepath%
  %
  % fill_do
  \brl@pgf_fill_do%
  % draw_do
  \brl@pgf_stroke_do%
}



\protected\def\brl@pgf_clipnext#1{%
  %
  \def\brl@pgf_clipnext_choose{#1}%
  \futurelet\brl@tmp\brl@pgf_clipnext_%
}


\def\brl@pgf_clipnext_{%
  \ifx\brl@tmp\brl@pgf_stroke%
    \def\brl@pgf_next{\brl@pgf_clipnext_stroke}%
  \else%
    \ifx\brl@tmp\brl@pgf_fillstroke%
      \def\brl@pgf_next{\brl@pgf_clipnext_fillstroke}%
    \else%
      \ifx\brl@tmp\brl@pgf_fill%
        \def\brl@pgf_next{\brl@pgf_clipnext_fill}%
      \else
        \def\brl@pgf_next{\brl@pgf_clipnext_clip}%
      \fi%
    \fi%
  \fi%
  \brl@pgf_next%
}

\def\brl@pgf_clipnext_stroke#1{%
  % stroke_on
  \brl@pgf@pointtoks{}%
  \let\brl@pgf_stroke_save_do\brl@pgf_stroke_save%
  % clip_on
  \brl@pgf@cliptoks{}%
  \let\brl@pgf_clip_save_do\brl@pgf_clip_save%
  \def\brl@pgf_clip@maxy{-\@M}%
  \def\brl@pgf_clip@miny{\@M}%
  %
  \brl@pgf_reset_localminmax%
  \the\brl@pgf@pathtoks%
  \brl@pgf@pathtoks{}%
  \brl@pgf_clip_save_closepath%
  % 
  % draw_do
  \brl@pgf_stroke_do%
  % clip_do
  \brl@pgf_clip_do%
}

\def\brl@pgf_clipnext_fill#1#2#3{%
  \brl@pgf_color{#2}%
  % stroke_off
  \let\brl@pgf_stroke_save_do\empty%
  % clip_on
  \brl@pgf@cliptoks{}%
  \let\brl@pgf_clip_save_do\brl@pgf_clip_save%
  \def\brl@pgf_clip@maxy{-\@M}%
  \def\brl@pgf_clip@miny{\@M}%
  %
  \brl@pgf_reset_localminmax%
  #3% eorule
  \the\brl@pgf@pathtoks%
  \brl@pgf@pathtoks{}%
  \brl@pgf_clip_save_closepath%
  %
  % fill_do
  \brl@pgf_fill_do%
  % clip_do
  \brl@pgf_clip_do%
}

\def\brl@pgf_clipnext_fillstroke#1#2#3{%
  \brl@pgf_color{#2}%
  % stroke_on
  \brl@pgf@pointtoks{}%
  \let\brl@pgf_stroke_save_do\brl@pgf_stroke_save%
  % clip_on
  \brl@pgf@cliptoks{}%
  \let\brl@pgf_clip_save_do\brl@pgf_clip_save%
  \def\brl@pgf_clip@maxy{-\@M}%
  \def\brl@pgf_clip@miny{\@M}%
  %
  \brl@pgf_reset_localminmax%
  #3% eorule
  \the\brl@pgf@pathtoks%
  \brl@pgf@pathtoks{}%
  \brl@pgf_clip_save_closepath%
  %
  % fill_do
  \brl@pgf_fill_do%
  % stroke_do
  \brl@pgf_stroke_do%
  % clip_do
  \brl@pgf_clip_do%
}

\def\brl@pgf_clipnext_clip{%
  % stroke_off
  \let\brl@pgf_stroke_save_do\empty%
  % clip_on
  \brl@pgf@cliptoks{}%
  \let\brl@pgf_clip_save_do\brl@pgf_clip_save%
  \def\brl@pgf_clip@maxy{-\@M}%
  \def\brl@pgf_clip@miny{\@M}%
  %
  \brl@pgf_reset_localminmax%
  \the\brl@pgf@pathtoks%
  \brl@pgf@pathtoks{}%
  \brl@pgf_clip_save_closepath%
  %
  % clip_do
  \brl@pgf_clip_do%
}




\def\brl@pgf_clip_save_closepath{%
  %
  % \brl@pgf@prevx, \brl@pgf@currx
  % \brl@pgf@prevy, \brl@pgf@curry
  % \brl@pgf@dx, \brl@pgf@dy
  % 
  \brl@pgf@x\pgf@xb%
  \brl@pgf@y\pgf@yb%
  \brl@pgf_transform{\brl@pgf@x}{\brl@pgf@y}%
  \brl@pgf@currx\numexpr\brl@pgf@x/\p@\relax%
  \brl@pgf@curry\numexpr\brl@pgf@y/\p@\relax%
  %
  \brl@pgf@dx\brl@pgf@prevx\advance\brl@pgf@dx-\brl@pgf@currx%
  \brl@pgf@dy\brl@pgf@prevy\advance\brl@pgf@dy-\brl@pgf@curry%
  \brl@pgf_clip_save_do%
  % 
  \brl@pgf@prevx\@M% initialize
  \brl@pgf@prevy\@M%
  \brl@pgf@prevxa\@M% initialize
  \brl@pgf@prevya\@M%
  \brl@pgf@prevxb\@M% initialize
  \brl@pgf@prevyb\@M%
  %
  %
  \pgf@xc\pgf@xb%
  \pgf@yc\pgf@yb%
}%


\def\brl@pgf_reset_localminmax{%
  \def\brl@pgf@localminx{\@M}%
  \def\brl@pgf@localminy{\@M}%
  \def\brl@pgf@localmaxx{-\@M}%
  \def\brl@pgf@localmaxy{-\@M}%
}


%%%%%%%%%%%%%%%%%%%%%%%%%%%%%%%%%%%%%%%%%%%%%%%%%%%%%%%%%%%%%%%%
%
% stroke, fill, clip
%



\def\brl@pgf_stroke_do{%
  \expandafter\brl@pgf_stroke_do_loop\the\brl@pgf@pointtoks\brl@end\brl@end\brl@end%
}

\def\brl@pgf_stroke_do_loop#1#2#3{%
  \def\brl@tmp{#1}%
  \ifx\brl@tmp\brl@@end%
  \else%
    \brl@pgf@currx#1\relax%
    \brl@pgf@curry#2\relax%
    \def\brl@pgf@pointsize{#3}% point size can be determined for each point.
    %
    \brl@pgf_drawpoint%
    \expandafter\brl@pgf_stroke_do_loop%
  \fi%
}

\def\brl@pgf_stroke_save{%
  \brl@pgf_stroke@pattern_do%
}%


% default
%
\def\brl@pgf_stroke@pattern_name_{%
  \brl@pgf_appendpoint%
}





%   step = 10
% 11 -> 20 -> 20
% 10 -> 19 -> 10
%  1 -> 10 -> 10
%  0 ->  9 ->  0
% -1       ->  0
% -9       ->  0
% -10      -> -10
%

\def\brl@pgf_fill_do{%
  %
  % 
  \brl@pgf_fill@y\brl@pgf_clip@miny\relax%
  %
  %
  % calculate \brl@pgf_fill@mody
  %
  \ifnum\brl@pgf_fill@y>\z@%
    \advance\brl@pgf_fill@y\brl@pgf_fill@stepym@ne\relax%
  \fi%
  \divide\brl@pgf_fill@y\brl@pgf_fill@stepy\relax%
  \brl@pgf_fill@mody\brl@pgf_fill@y%
  \count@\brl@pgf_fill@y%
  \multiply\brl@pgf_fill@y\brl@pgf_fill@stepy\relax%
  % 
  % 
  %
  % step = 10
  % 11 -> 1
  % 10 -> 0
  %  9 -> 9
  %  0 -> 0
  % -1 -> -1
  % -9 -> -9
  % -10 -> 0
  % -11 -> -1
  \divide\count@\brl@pgf_fill@maxmody\relax%
  \multiply\count@\brl@pgf_fill@maxmody\relax%
  \advance\brl@pgf_fill@mody-\count@% 
  %
  \ifnum\brl@pgf_fill@mody<\z@%
    \advance\brl@pgf_fill@mody\brl@pgf_fill@maxmody\relax%
  \fi%
  %
  %
  \brl@pgf_fill_do_loop%
}


\def\brl@pgf_fill_do_loop{%
  % 
  \ifnum\brl@pgf_clip@maxy>\brl@pgf_fill@y% data of maxy is not defined.
    % 
    %
    \brl@pgf@curry\brl@pgf_fill@y%
    %
    %
    %
    % 
    \brl@pgf_fill@num\@ne%
    \brl@pgf@settoks{}%
    \def\brl@pgf_fill_clip_next{\brl@pgf_fill_clip_y}%
    \expandafter\brl@pgf_fill_clip_y\the\brl@pgf@clipalltoks\brl@end\brl@end\brl@end\brl@end%
    % 
    \advance\brl@pgf_fill@y\brl@pgf_fill@stepy\relax%
    % 
    % 
    \advance\brl@pgf_fill@mody\@ne%
    \ifnum\brl@pgf_fill@maxmody=\brl@pgf_fill@mody%
      \brl@pgf_fill@mody\z@%
    \fi%
    % 
    %
    \expandafter\brl@pgf_fill_do_loop%
  \fi%
}



%   x -> x/10*10
%  14 -> 10
%   4 -> 0
%  -4 -> 0
% -14 -> -10

%
% #1, #2 : xpos
%
% #1 = #2 may happen.
%
\def\brl@pgf_fill_do_dots#1#2{% 
  \def\brl@pgf_fill@startx{#1}%
  \ifx\brl@pgf_fill@startx\brl@@end%
  \else%
    % 
    %
    % \brl@pgf_fill@x mod \brl@pgf_fill@stepx + \brl@pgf_fill@step_\mody
    %
    % step = 10
    % 11 -> 10
    % 10 -> 10
    % 9 -> 0
    % 0 -> 0
    % -1 -> -10
    % -9 -> -10
    % -10 -> -20
    % -11 -> -20
    %
    \brl@pgf_fill@x#1\relax%
    \ifnum\brl@pgf_fill@x<\z@%
      \advance\brl@pgf_fill@x-\brl@pgf_fill@stepx\relax%
    \fi%
    \divide\brl@pgf_fill@x\brl@pgf_fill@stepx\relax%
    \multiply\brl@pgf_fill@x\brl@pgf_fill@stepx\relax%
    %
    \csname brl@pgf_fill@pattern_\brl@pgf_fill@pattern_name _\the\brl@pgf_fill@mody _0\endcsname% init
    %
    %
    \brl@pgf_fill_do_search_startx%
    %
    \def\brl@pgf_fill@maxx{#2}%
    %
    \brl@pgf_fill_do_dots_loop%
    % 
    \expandafter\brl@pgf_fill_do_dots%
  \fi%
}

\def\brl@pgf_fill_do_search_startx{%
  \ifnum\brl@pgf_fill@startx>\brl@pgf_fill@x%
    \brl@pgf_fill@dostepx%
    \expandafter\brl@pgf_fill_do_search_startx%
  \fi%
}


\def\brl@pgf_fill_do_dots_loop{%
  \ifnum\brl@pgf_fill@maxx>\brl@pgf_fill@x%
    \brl@pgf@currx\brl@pgf_fill@x%
    %
    \brl@pgf_drawpoint_noclip%
    % 
    \brl@pgf_fill@dostepx%
    \expandafter\brl@pgf_fill_do_dots_loop%
  \fi%
}





\def\brl@pgf_fill_clip_y#1#2#3#4{% #1: eorule #2: ymin, #3: ymax, #4: path
  \def\brl@tmp{#1}%
  % 
  \ifx\brl@tmp\brl@@end%
    % OK
    %
    \brl@pgf@clipytoks{}%
    %
    \expandafter\brl@pgf_clip_truncate\the\brl@pgf@cliptoks\brl@end\brl@end\brl@end\brl@end\brl@end%
    {%
      \let\brl@qs_choose_def\brl@qs_choose_def_set_p% union of disjoint sets of intervals
      \expandafter\brl@qs\expandafter{\the\brl@pgf@clipytoks}%
    }% quick sort
    %
    \expandafter\brl@pgf_appendsettoks\expandafter{\brl@qs@res}%
    {%
      \let\brl@qs_choose\brl@qs_choose_set% intersection of union of disjoint sets of intervals
      \expandafter\brl@qs\expandafter{\the\brl@pgf@settoks}%
    }% quick sort
    \expandafter\brl@pgf_fill_do_dots\brl@qs@res\brl@end\brl@end%
    % 
  \else%
    \ifnum#2<\brl@pgf@curry%
      \ifnum#3>\brl@pgf@curry% inside
        % 
        % 
        \brl@pgf@clipytoks{}%
        \brl@pgf_clip_truncate#4\brl@end\brl@end\brl@end\brl@end\brl@end%
        {%
          #1% eorule
          \let\brl@qs_choose_def\brl@qs_choose_def_set_p% union of disjoint sets of intervals
          \expandafter\brl@qs\expandafter{\the\brl@pgf@clipytoks}%
        }% quick sort
        %
        \advance\brl@pgf_fill@num\@ne%
        \expandafter\brl@pgf_appendsettoks\expandafter{\brl@qs@res}%
        %
      \else%
        \def\brl@pgf_fill_clip_next{\brl@pgf_fill_clip_end}%
      \fi%
    \else%
      \def\brl@pgf_fill_clip_next{\brl@pgf_fill_clip_end}%
    \fi%
    \expandafter\brl@pgf_fill_clip_next%
  \fi%
}


\def\brl@pgf_fill_clip_end#1\brl@end\brl@end\brl@end\brl@end{}%



%%%%%%%%%%%%%%%%%%%%%%%%%%%%%%%%%%%%%%%%%%%%%%%%%%%%%%%%%%%%%
%
% \brl@pgf_clip_save
%
% make clip path
%
% input :
% \brl@pgf@dx, \brl@pgf@dy
% \brl@pgf@currx, \brl@pgf@curry
% \brl@pgf@prevx, \brl@pgf@prevy
%
% output : \brl@pgf@cliptoks
% \brl@pgf_clip@miny
% \brl@pgf_clip@maxy
%

\def\brl@pgf_clip_save{%
  \ifnum\brl@pgf@prevx=\@M%
  \else%
    \ifnum\brl@pgf@dy>\z@% \brl@pgf@prevy > \brl@pgf@curry
      %
      \brl@pgf_clip_save_slope{\brl@pgf@currx}{\brl@pgf@curry}{\brl@pgf@prevy}\@ne%
    \else%
      \ifnum\brl@pgf@dy<\z@% \brl@pgf@prevy < \brl@pgf@curry
        %
        \brl@pgf_clip_save_slope{\brl@pgf@prevx}{\brl@pgf@prevy}{\brl@pgf@curry}\m@ne%
      \fi%
      % ignore horizontal lines
    \fi%
  \fi%
}

\def\brl@pgf_clip_save_slope#1#2#3#4{% (#1,#2) (??,#3) : #2<#3
  \edef\brl@tmp{{{\the#1}{\the#2}{\the#3}{\the\brl@pgf@dx/\the\brl@pgf@dy}#4}}%
  \expandafter\brl@pgf_appendcliptoks\brl@tmp%
  \ifnum\brl@pgf_clip@miny>#2%
    \edef\brl@pgf_clip@miny{\the#2}%
  \fi%
  \ifnum\brl@pgf_clip@maxy<#3%
    \edef\brl@pgf_clip@maxy{\the#3}%
  \fi%
}

%

%%%%%%%%%%%%%%%%%%%%%%%%%%%%%%%%%%%%%%%%%%%%%%%%%%%%%%%%%%%%%%%%
%
% \brl@pgf_clip_truncate
%
%
% input:
% \brl@pgf@cliptoks, \brl@pgf@curry
% 
% outout: x coordinates width \@ne (upwards) or \m@ne (downwards)
% 
%
% algorithm : consider the intersecting points of
% the line (\brl@pgf@prevx, \brl@pgf@prevy) to (\brl@pgf@currx,\brl@pgf@curry) 
%

%
%
%   step = 10
% 11 -> 20 -> 20
% 10 -> 19 -> 10
%  1 -> 10 -> 10
%  0 ->  9 ->  0
% -1       ->  0
% -9       ->  0
% -10      -> -10

\def\brl@pgf_clip_truncate#1#2#3#4#5{% #4 = \brl@pgf@dx/\brl@pgf@dy, #5 = \@ne or \m@ne
  \def\brl@tmp{#1}%
  \ifx\brl@tmp\brl@@end%
  \else%
    %
    \ifnum#2>\brl@pgf@curry% #2 <= \brl@pgf@curry < #3
    \else%
      \ifnum#3>\brl@pgf@curry%
        \edef\brl@tmp{{{\the\numexpr(\brl@pgf@curry-#2)*#4+#1\relax}#5}}%
        \expandafter\brl@pgf_appendclipytoks\brl@tmp%
        %
      \fi%
    \fi%
    \expandafter\brl@pgf_clip_truncate%
  \fi%
}

%%%%%%%%%%%%%%%%%%%%%%%%%%%%%%%%%%%%%%%%%%%%%%%%%%%%%%%%%%%%%%%%
%
% input : \brl@pgf_clip@miny, \brl@pgf_clip@maxy, \brl@pgf@cliptoks
%
\def\brl@pgf_clip_do{%
  \edef\brl@tmp{{{\brl@pgf_clipnext_choose}{\brl@pgf_clip@miny}{\brl@pgf_clip@maxy}{\the\brl@pgf@cliptoks}}}%
  \expandafter\brl@pgf_appendclipalltoks\brl@tmp%
}


%%%%%%%%%%%%%%%%%%%%%%%%%%%%%%%%%%%%%%%%%%%%%%%%%%%%%%%%%%%%%%%%
%
% transformation
%

% Transformation matrix


\def\brl@pgf@pt_aa{1.0}\def\brl@pgf@pt_ab{0.0}
\def\brl@pgf@pt_ba{0.0}\def\brl@pgf@pt_bb{1.0}

\protected\def\brl@pgf_transform#1#2{%
  \pgf@pt@temp#1%
  #1\dimexpr\brl@pgf@pt_aa\pgf@pt@temp+\brl@pgf@pt_ba#2+\brl@pgf@pt_x\relax%
  #2\dimexpr\brl@pgf@pt_ab\pgf@pt@temp+\brl@pgf@pt_bb#2+\brl@pgf@pt_y\relax%
}

\protected\def\brl@pgf_transformcm#1#2#3#4#5#6{%
  %
  \brl@pgf@xc#5\relax%
  \brl@pgf@yc#6\relax%
  % 
  \brl@pgf@x#1\p@%
  \brl@pgf@y#2\p@%
  % 
  \brl@pgf@xa\brl@pgf@pt_aa\brl@pgf@x\advance\brl@pgf@xa\brl@pgf@pt_ba\brl@pgf@y%
  \brl@pgf@ya\brl@pgf@pt_ab\brl@pgf@x\advance\brl@pgf@ya\brl@pgf@pt_bb\brl@pgf@y%
  % 
  \brl@pgf@x#3\p@%
  \brl@pgf@y#4\p@%
  % 
  \brl@pgf@xb\brl@pgf@pt_aa\brl@pgf@x\advance\brl@pgf@xb\brl@pgf@pt_ba\brl@pgf@y%
  \brl@pgf@yb\brl@pgf@pt_ab\brl@pgf@x\advance\brl@pgf@yb\brl@pgf@pt_bb\brl@pgf@y%
  % 
  \advance\brl@pgf@pt_x\brl@pgf@pt_aa\brl@pgf@xc\advance\brl@pgf@pt_x\brl@pgf@pt_ba\brl@pgf@yc%
  \advance\brl@pgf@pt_y\brl@pgf@pt_ab\brl@pgf@xc\advance\brl@pgf@pt_y\brl@pgf@pt_bb\brl@pgf@yc%
  % 
  \edef\brl@pgf@pt_aa{\pgf@sys@tonumber{\brl@pgf@xa}}%
  \edef\brl@pgf@pt_ab{\pgf@sys@tonumber{\brl@pgf@ya}}%
  \edef\brl@pgf@pt_ba{\pgf@sys@tonumber{\brl@pgf@xb}}%
  \edef\brl@pgf@pt_bb{\pgf@sys@tonumber{\brl@pgf@yb}}%
  %
}



\protected\def\brl@pgf_setdash#1#2{% #1 : pattern, #2 : phase
  %
  \ifcsname brl@pgf_stroke@pattern_name_#1\endcsname%
  \else% if the cs is not defined, then create it and register it.
    \global\advance\brl@pgf_stroke@pattern_num\@ne%
    % 
    \def\brl@pgf_stroke@pattern_point{\brl@pgf_appendpoint}%
    \count@\z@%
    \expandafter\brl@pgf_setdash_pattern#1,\brl@end,%
    % the last = 0
    \global\expandafter\let\csname brl@pgf_stroke@pattern_name_#1\expandafter\endcsname\csname brl@pgf_stroke@pattern_\the\brl@pgf_stroke@pattern_num _0\endcsname%
    \global\expandafter\let\csname brl@pgf_stroke@pattern_\the\brl@pgf_stroke@pattern_num _\the\count@\expandafter\endcsname\csname brl@pgf_stroke@pattern_\the\brl@pgf_stroke@pattern_num _0\endcsname%
    % 
  \fi%
  \expandafter\let\expandafter\brl@pgf_stroke@pattern_do\csname brl@pgf_stroke@pattern_name_#1\endcsname%
}


\def\brl@pgf_setdash_pattern#1,{%
  \def\brl@tmp{#1}%
  \ifx\brl@tmp\brl@@end%
  \else%
    % 
    \brl@pgf_stroke@pattern_phase\numexpr\dimexpr#1\relax/\p@\relax%
    \ifnum\brl@pgf_stroke@pattern_phase<\@ne%
      \brl@pgf_stroke@pattern_phase\@ne%
    \fi%
    \brl@pgf_setdash_loop%
    \ifx\brl@pgf_stroke@pattern_point\empty%
      \def\brl@pgf_stroke@pattern_point{\brl@pgf_appendpoint}% on
    \else%
      \let\brl@pgf_stroke@pattern_point\empty% off
    \fi%
    \expandafter\brl@pgf_setdash_pattern%
  \fi%
}

\def\brl@pgf_setdash_loop{%
  \edef\brl@tmp{\the\count@}%
  \advance\count@\@ne%
  \expandafter\xdef\csname brl@pgf_stroke@pattern_\the\brl@pgf_stroke@pattern_num _\brl@tmp\endcsname{%
    \let\noexpand\brl@pgf_stroke@pattern_do\expandafter\noexpand\csname brl@pgf_stroke@pattern_\the\brl@pgf_stroke@pattern_num _\the\count@\endcsname%
    \brl@pgf_stroke@pattern_point%
  }%
  %
  \ifnum\brl@pgf_stroke@pattern_phase>\@ne%
    \advance\brl@pgf_stroke@pattern_phase\m@ne%
    \expandafter\brl@pgf_setdash_loop%
  \fi%
}


\protected\def\brl@pgf_appendpoint_single{%
  \edef\brl@tmp{{{\the\brl@pgf@currx}{\the\brl@pgf@curry}{\brl@pgf@pointsize}}}% currx, curry, pointsize
  \expandafter\brl@pgf_appendpointtoks\brl@tmp%
}%

\protected\def\brl@pgf_appendpoint_double{%
  \edef\brl@tmp{{{\the\brl@pgf@currxa}{\the\brl@pgf@currya}{\brl@pgf@pointsize@p}%
      {\the\brl@pgf@currxb}{\the\brl@pgf@curryb}{\brl@pgf@pointsize@p}}}% currx, curry, pointsize
  \expandafter\brl@pgf_appendpointtoks\brl@tmp%
}%

\protected\def\brl@pgf_appendpoint_triple{%
  \edef\brl@tmp{{{\the\brl@pgf@currx}{\the\brl@pgf@curry}{\brl@pgf@pointsize@p}%
      {\the\brl@pgf@currxa}{\the\brl@pgf@currya}{\brl@pgf@pointsize@p}%
      {\the\brl@pgf@currxb}{\the\brl@pgf@curryb}{\brl@pgf@pointsize@p}}}% currx, curry, pointsize
  \expandafter\brl@pgf_appendpointtoks\brl@tmp%
}%



% Line width

% ultra thin 0.1pt
% very thin 0.2pt
% thin 0.4pt
% semithick 0.6pt
% thick 0.8pt
% very thick 1.2pt
% ultra thick 1.6pt


\protected\def\brl@pgf_setlinewidth#1{%
  \@tempdima#1\relax%
  \ifdim\@tempdima<1.5\p@%
    \brl@pgf_setlinewidth_single%
  \else%
    \let\brl@pgf_setpoint_line\brl@pgf_setpoint_line_double%
    \let\brl@pgf_lineto_do_normalize\brl@pgf_lineto_do_normalize_double%
    \let\brl@pgf_setpoint_curve\brl@pgf_setpoint_curve_double%
    \let\brl@pgf_curveto_do_normalize\brl@pgf_curveto_do_normalize_double%
    \ifdim\@tempdima<3.0\p@%
      \brl@pgf_setlinewidth_double%
    \else%
      \brl@pgf_setlinewidth_triple%
    \fi%
  \fi%
}%


\def\brl@pgf_setlinewidth_single{%
  \let\brl@pgf_setpoint_line\brl@pgf_setpoint_line_single%  
  \let\brl@pgf_lineto_do_normalize\empty%
  \let\brl@pgf_setpoint_curve\brl@pgf_setpoint_curve_single%  
  \let\brl@pgf_curveto_do_normalize\brl@pgf_normalize_novec%
  \let\brl@pgf_appendpoint\brl@pgf_appendpoint_single%
  % 
  \ifdim\@tempdima<0.5\p@%
    \def\brl@pgf@pointsize{1}%
  \else%
    \ifdim\@tempdima<\p@%
      \def\brl@pgf@pointsize{2}%
    \else%
      \def\brl@pgf@pointsize{3}%
    \fi%
  \fi%
}%

\def\brl@pgf_setlinewidth_double{%
  \let\brl@pgf_appendpoint\brl@pgf_appendpoint_double%
  %
  \ifdim\@tempdima<2.0\p@%
    \def\brl@pgf@pointsize{4}%
    \def\brl@pgf@pointsize@p{1}% pointsize
    \def\brl@pgf@pointsize@d{2.5}% distance of two points
  \else%
    \ifdim\@tempdima<2.5\p@%
      \def\brl@pgf@pointsize{5}%
      \def\brl@pgf@pointsize@p{2}%
      \def\brl@pgf@pointsize@d{2.75}%
    \else%
      \def\brl@pgf@pointsize{6}%
      \def\brl@pgf@pointsize@p{3}%
      \def\brl@pgf@pointsize@d{3}%
    \fi%
  \fi%
}

\def\brl@pgf_setlinewidth_triple{%
  \let\brl@pgf_appendpoint\brl@pgf_appendpoint_triple%
  %
  \ifdim\@tempdima<3.5\p@%
    \def\brl@pgf@pointsize{7}%
    \def\brl@pgf@pointsize@p{1}% pointsize
    \def\brl@pgf@pointsize@d{4}% distance of two points
  \else%
    \ifdim\@tempdima<4.0\p@%
      \def\brl@pgf@pointsize{8}%
      \def\brl@pgf@pointsize@p{2}%
      \def\brl@pgf@pointsize@d{5}%
    \else%
      \def\brl@pgf@pointsize{9}%
      \def\brl@pgf@pointsize@p{3}%
      \def\brl@pgf@pointsize@d{6}%
    \fi%
  \fi%
}

%% if the option includes `double', then, double is cancelled and the line is doubled.

\protected\def\brl@pgf_setlinewidth_setdouble#1{%
  \@tempdima#1\relax%
  \let\brl@pgf_setpoint_line\brl@pgf_setpoint_line_double%
  \let\brl@pgf_lineto_do_normalize\brl@pgf_lineto_do_normalize_double%
  \let\brl@pgf_setpoint_curve\brl@pgf_setpoint_curve_double%
  \let\brl@pgf_curveto_do_normalize\brl@pgf_curveto_do_normalize_double%
  \brl@pgf_setlinewidth_double%
}%

% input :
% \brl@pgf@currx, \brl@pgf@curry, \brl@pgf@pointsize
% \brl@pgf@clipalltoks (clip path)
%
% output :
% \brl@pgf@\brl@pgfminx, \brl@pgf@miny, \brl@pgf@maxx, \brl@pgf@maxy
% \brl@pgf@minxaty_y
% \brl@pgf@maxxaty_y
% \brl@pgf@size_x_y


\def\brl@pgf_drawpoint{%
  %
  \def\brl@pgf_drawpoint_clip_next{\brl@pgf_drawpoint_clip}%
  % 
  \expandafter\brl@pgf_drawpoint_clip\the\brl@pgf@clipalltoks\brl@end\brl@end\brl@end\brl@end%
}

\def\brl@pgf_drawpoint_clip#1#2#3#4{% #1: eorule #2: ymin, #3: ymax, #4: path
  \def\brl@tmp{#1}%
  % 
  \ifx\brl@tmp\brl@@end%
    % OK
    %
    \expandafter\brl@pgf_drawpoint_noclip%
    % 
  \else%
    \ifnum#2<\brl@pgf@curry%
      \ifnum#3>\brl@pgf@curry% inside
        % 
        % 
        \brl@pgf@clipytoks{}%
        \brl@pgf_clip_truncate#4\brl@end\brl@end\brl@end\brl@end\brl@end%
        {%
          #1% eorule
          \expandafter\brl@qs\expandafter{\the\brl@pgf@clipytoks}%
        }% quick sort
        %
        \def\brl@pgf_drawpoint_clip__next{\brl@pgf_drawpoint_clip_}%
        \expandafter\brl@pgf_drawpoint_clip_\brl@qs@res\brl@end\brl@end%
      \else%
        %
        \def\brl@pgf_drawpoint_clip_next{\brl@pgf_drawpoint_clip_end}%
      \fi%
    \else%
      %
      \def\brl@pgf_drawpoint_clip_next{\brl@pgf_drawpoint_clip_end}%
    \fi%
    \expandafter\brl@pgf_drawpoint_clip_next%
  \fi%
}

\def\brl@pgf_drawpoint_clip_end#1\brl@end\brl@end\brl@end\brl@end{}%

\def\brl@pgf_drawpoint_clip_#1#2{%
  \def\brl@tmp{#1}%
  \ifx\brl@tmp\brl@@end%
    % larger than all.
    \def\brl@pgf_drawpoint_clip_next{\brl@pgf_drawpoint_clip_end}%
  \else%
    \ifnum#1<\brl@pgf@currx%
      \ifnum#2>\brl@pgf@currx%
        % OK
        \def\brl@pgf_drawpoint_clip__next{\brl@pgf_drawpoint_clip__end}%
        % \else
        % not yet determined
      \fi%
    \else%
      \def\brl@pgf_drawpoint_clip_next{\brl@pgf_drawpoint_clip_end}%
      %
      \def\brl@pgf_drawpoint_clip__next{\brl@pgf_drawpoint_clip__end}%
    \fi%
    \expandafter\brl@pgf_drawpoint_clip__next%
  \fi%
}

\def\brl@pgf_drawpoint_clip__end#1\brl@end\brl@end{}%



\def\brl@pgf_drawpoint_noclip{%
  \brl@pgf_drawpoint_init%
  % \brl@pgf@currx, \brl@pgf@curry -> \brl@pgf@transx, \brl@pgf@transy 
  %
  \ifnum\brl@pgf@pointsize>\z@%
    %
    % set global min, max
    %
    \ifnum\brl@pgf@transx<\brl@pgf@minx\relax%
      \global\let\brl@pgf@minx\brl@pgf@transx%
    \fi%
    \ifnum\brl@pgf@transy<\brl@pgf@miny\relax%
      \global\let\brl@pgf@miny\brl@pgf@transy%
    \fi%
    \ifnum\brl@pgf@transx>\brl@pgf@maxx\relax%
      \global\let\brl@pgf@maxx\brl@pgf@transx%
    \fi%
    \ifnum\brl@pgf@transy>\brl@pgf@maxy\relax%
      \global\let\brl@pgf@maxy\brl@pgf@transy%
    \fi%
    % 
    % set local min, max
    %
    \edef\brl@tmpmin{brl@pgf@minxaty_\brl@pgf@transy}%
    \edef\brl@tmpmax{brl@pgf@maxxaty_\brl@pgf@transy}%
    \ifcsname\brl@tmpmin\endcsname%
      \ifnum\csname\brl@tmpmin\endcsname>\brl@pgf@transx\relax%
        \global\expandafter\let\csname\brl@tmpmin\endcsname\brl@pgf@transx%
      \fi%
      \ifnum\csname\brl@tmpmax\endcsname<\brl@pgf@transx\relax%
        \global\expandafter\let\csname\brl@tmpmax\endcsname\brl@pgf@transx%
      \fi%
    \else%
      \global\expandafter\let\csname\brl@tmpmin\endcsname\brl@pgf@transx%
      \global\expandafter\let\csname\brl@tmpmax\endcsname\brl@pgf@transx%
    \fi%
    % 
    \global\expandafter\let\csname brl@pgf@size_\brl@pgf@transx _\brl@pgf@transy\endcsname\brl@pgf@pointsize% set
    % 
  \else%
    \global\expandafter\let\csname brl@pgf@size_\brl@pgf@transx _\brl@pgf@transy\endcsname\@undefined% unset
  \fi%
}


%%%%%%%%%%%%%%%%%%%%%%%%%%%%%%%%%%%%%%%%%%%%%%%%%%%%%%%%%%%%%%%%
%
% text
%

%
%

\protected\def\brl@pgf_put#1#2#3{%
  \brl@pgf@x#1\relax%
  \brl@pgf@y#2\relax%
  \brl@pgf_transform{\brl@pgf@x}{\brl@pgf@y}%
  \brl@pgf@ptx\numexpr\brl@pgf@x/\p@\relax%
  \brl@pgf@pty\numexpr\brl@pgf@y/\p@\relax%
  \brl@pgf_put_loop#3\brl@end%
}

\def\brl@pgf_put_loop#1{%
  \def\brl@tmp{#1}%
  \ifx\brl@tmp\brl@@end%
  \else%
    \edef\brl@tmp{{\brl@pgf_hbox{\the\brl@pgf@ptx}{\the\brl@pgf@pty}\unexpanded{{#1}}}}%
    \expandafter\brl@pgf_appendnodetoks\brl@tmp% save position
    %
    % 
    %
    \advance\brl@pgf@pty-\brl@pgf@oneletter\relax% advance one letter.
    \expandafter\brl@pgf_put_loop%
  \fi%
}


%
% color -> fill pattern
%
% #1 : gray_0.1, rgb_0,0,1 etc.
%
% if the corresponding pattern is defined, then use it.
%
% Otherwise, use the grayscale.
%
% \brl@pgf_fill@pattern_name : data of fill pattern
% \brl@pgf_fill@pattern_name_num : patterns
%



% \numexpr / \relax
% 19/10  = 1
% 10/10  = 1
% 5/10   = 1
% 4/10   = 0
% -4/10  = 0
% -5/10  = -1
% -9/10  = -1
% -10/10 = -1
% -14/10 = -1
% -15/10 = -2
% -19/10 = -2



% \divide
% 19/10  = 1   -> 20
% 10/10  = 1   -> 10
% 9/10   = 0   -> 10
% 4/10   = 0   -> 10
% 1/10   = 0   -> 10
% 0/10   = 0   -> 0
% -9/10  = 0   -> 0
% -10/10 = -1  -> -10
% -19/10 = -1  -> -10
% -20/10 = -2  -> -20

%%%%%%%%%%%%%%%%%%%%%%%%%%%%%%%%%%%%%%%%%%%%%%%%%%%%%%%%%%%%%
%
% \brl@pgf_filldef
%
% define fill pattern
%

\def\brl@pgf_filldef_getcolor#1#2#3#4#5{#4_#5}



% \brl@pgf_filldef{10}{1}{1}{%
%   {0101}%
%   {1010}%
%   {0101}%
%   {1010}%
% }

%
% #1 : pattern_name
% #2 : xstep (the distance between two dots)
% #3 : ystep (the distance between two dots)
% #4 : set of pattern


\def\brl@pgf_filldef_save#1#2#3#4{%
  \expandafter\gdef\csname brl@pgf_filldef_save@#1\endcsname{%
    \brl@pgf_filldef{#1}{#2}{#3}{#4}%
  }%
}



\def\brl@pgf_plot_filldef#1#2#3#4{%
  \edef\brl@pgf_filldef@pattern_name{#1}%
  \brl@pgf_filldef@curry\z@%
  \def\brl@pgf_filldef@maxx{-1}%
  \edef\brl@pgf_filldef@minstepx{\the\numexpr\brl@pgf_filldef@plotunit*#2\relax}%
  \edef\brl@pgf_filldef@minstepy{\the\numexpr\brl@pgf_filldef@plotunit*#3\relax}%
  %
  \brl@pgf_plot_filldef_count#4\brl@end% to reverse the order, count the number of patterns
  \edef\brl@pgf_filldef@maxy{\the\brl@pgf_filldef@curry}%
  %
  % plot
  %
  \brl@pgf_plot_filldef_pattern#4%
  % 
  \edef\brl@tmp{%
    \gdef\noexpand\brl@pgf_fill@stepx{\brl@pgf_filldef@maxx}%
    \gdef\noexpand\brl@pgf_fill@stepy{\brl@pgf_filldef@minstepy}%
    \gdef\noexpand\brl@pgf_fill@stepym@ne{\the\numexpr\brl@pgf_filldef@minstepy-\@ne\relax}%
    \gdef\noexpand\brl@pgf_fill@maxmody{\brl@pgf_filldef@maxy}%
    \gdef\noexpand\brl@pgf_fill@pattern_name{\brl@pgf_filldef@pattern_name}%
  }%
  \global\expandafter\let\csname brl@pgf_fill@pattern_\brl@pgf_filldef@pattern_name\endcsname\brl@tmp%
  %
  \brl@tmp% execute right now.
}

\def\brl@pgf_plot_filldef_count#1{%
  \def\brl@tmp{#1}%
  \ifx\brl@tmp\brl@@end%
  \else%
    \advance\brl@pgf_filldef@curry\@ne%
    \expandafter\brl@pgf_plot_filldef_count%
  \fi%
}


\def\brl@pgf_plot_filldef_pattern#1{%
  %
  %
  %
  \advance\brl@pgf_filldef@curry\m@ne%
  \brl@pgf_filldef@currx\z@%
  \brl@pgf_filldef@dotnum\z@%
  \brl@pgf_filldef@prevx\z@%
  \def\brl@pgf_filldef@prevsize{0}%
  \edef\brl@pgf_filldef@defname{brl@pgf_fill@pattern_\brl@pgf_filldef@pattern_name _\the\brl@pgf_filldef@curry _}% for fixed pattern and fixed y
  % 
  \brl@pgf_plot_filldef_pattern_loop#1\brl@end%
  \ifnum\brl@pgf_filldef@curry>\z@%
    \expandafter\brl@pgf_plot_filldef_pattern%
  \fi%
}



% 0001100010
%
% _0 \advance 3 \let\brl@pgf_fill@dostepx _1
% _1 \advance 1 \let\brl@pgf_fill@dostepx _2
% _2 \advance 4 \let\brl@pgf_fill@dostepx _3
% _3 \advance 5 \let\brl@pgf_fill@dostepx _1
%
% 0000100000
% _0 \advance 4 \let\brl@pgf_fill@dostepx _1
% _1 \advance 10
%
% 1000000000
% _0 \let\brl@pgf_fill@dostepx _1
% _1 \advance 10
%
% 0000000000
% _0 \advance 10000



\def\brl@pgf_plot_filldef_pattern_loop#1{%
  \def\brl@pgf_filldef@currsize{#1}%
  \ifx\brl@pgf_filldef@currsize\brl@@end%
    %
    % end
    \expandafter\brl@pgf_plot_filldef_pattern_end%
    % 
  \else%
    %
    % loop
    %
    \ifnum\brl@pgf_filldef@currsize>\z@%
      %
      %
%      \ifnum\brl@pgf_filldef@dotnum=\z@%
      \ifcase\brl@pgf_filldef@dotnum%
        \edef\brl@pgf_filldef@firstx{\the\brl@pgf_filldef@currx}%
        \let\brl@pgf_filldef@firstsize\brl@pgf_filldef@currsize%
      \fi%
      % 
      \edef\brl@tmp{\the\numexpr\brl@pgf_filldef@currx-\brl@pgf_filldef@prevx\relax}%
      \expandafter\xdef\csname\brl@pgf_filldef@defname\the\brl@pgf_filldef@dotnum\endcsname{%
        %
        \unless\ifnum\brl@pgf_filldef@currsize=\brl@pgf_filldef@prevsize% the same size
          \def\noexpand\brl@pgf@pointsize{\brl@pgf_filldef@currsize}%
          % 
        \fi%
        \ifnum\brl@tmp>\z@% this happens when the pattern is 1.... 
          \advance\brl@pgf_fill@x\brl@tmp\relax%
        \fi%
        \let\noexpand\brl@pgf_fill@dostepx%
        \expandafter\noexpand\csname\brl@pgf_filldef@defname\the\numexpr\brl@pgf_filldef@dotnum+1\relax\endcsname%
      }%
      %
      %
      \let\brl@pgf_filldef@prevsize\brl@pgf_filldef@currsize%
      % 
      \brl@pgf_filldef@prevx\brl@pgf_filldef@currx%
      \advance\brl@pgf_filldef@dotnum\@ne%
    \fi%
    %
    %
    \advance\brl@pgf_filldef@currx\brl@pgf_filldef@minstepx\relax%
    % 
    \ifnum\brl@pgf_filldef@maxx<\brl@pgf_filldef@currx%
      \edef\brl@pgf_filldef@maxx{\the\brl@pgf_filldef@currx}%
    \fi%
    %
    %
    %
    \expandafter\brl@pgf_plot_filldef_pattern_loop%
  \fi%
}


\def\brl@pgf_plot_filldef_pattern_end{%
  %
  \ifcase\brl@pgf_filldef@dotnum%
    % 
    \expandafter\gdef\csname\brl@pgf_filldef@defname 0\endcsname{\brl@pgf_fill@x\@M}%
  \or%
    \expandafter\xdef\csname\brl@pgf_filldef@defname 1\endcsname{\advance\brl@pgf_fill@x\brl@pgf_filldef@maxx\relax}%
  \else%
    %
    %
    \expandafter\xdef\csname\brl@pgf_filldef@defname\the\brl@pgf_filldef@dotnum\endcsname{%
      %
      \unless\ifnum\brl@pgf_filldef@prevsize=\brl@pgf_filldef@firstsize% the same size
        \def\noexpand\brl@pgf@pointsize{\brl@pgf_filldef@firstsize}%
        % 
      \fi%
      \advance\brl@pgf_fill@x%
      \the\numexpr\brl@pgf_filldef@maxx-\brl@pgf_filldef@prevx+\brl@pgf_filldef@firstx\relax\relax%
      \let\noexpand\brl@pgf_fill@dostepx\expandafter\noexpand\csname\brl@pgf_filldef@defname 1\endcsname%
    }% add loop
  \fi%
} 




\def\brl@pgf_color#1{%
  {%
    \def\brl@tmp{#1}%
    % 
    \ifx\brl@tmp\empty%
      \def\brl@tmp{black}%
    \fi%
    \brl@orig@color{\brl@tmp}% gives \XC@current@color -> \xcolor@ {}{rgb 1 0 0}{rgb}{1,0,0}
    %
    \edef\brl@tmp{\expandafter\brl@pgf_filldef_getcolor\XC@current@color}%
    % 
    \ifcsname brl@pgf_fill@pattern_\brl@tmp\endcsname%
      \csname brl@pgf_fill@pattern_\brl@tmp\endcsname%
    \else%
      \ifcsname brl@pgf_filldef_save@\brl@tmp\endcsname% user definition
        \csname brl@pgf_filldef_save@\brl@tmp\endcsname% exec user definition
      \else%
        \expandafter\brl@pgf_gray\brl@tmp\brl@end%
        % 
        \ifcsname brl@pgf_fill@pattern_\the\count@\endcsname%
          \csname brl@pgf_fill@pattern_\the\count@\endcsname%
        \else%
          \csname brl@pgf_filldef_save@\the\count@\endcsname% exec predefinition
        \fi%
      \fi%
    \fi%
  }%
}



\def\brl@pgf_gray#1_{%
  \csname brl@pgf_model_#1_\endcsname%
}

\def\brl@pgf_init{% init for plot
  \brl@pgf@clipalltoks{}%
  \brl@pgf@pathtoks{}%
  \global\brl@pgf@invoketoks{}%
}


%%
%% matrix
%%
%%
\def\brl@pgf_plot_matrix_init{%
  \let\brl@pgf@matrix@endcell_sub\brl@pgf@matrix@endcell_plot%
  \let\brl@pgf@matrix@step@column_sub\brl@pgf@matrix@step@column_plot%
  \let\brl@pgf_matrix_sub\brl@pgf_matrix_plot%
}


\def\brl@pgf@matrix@step@column_plot{%
  \let\pgfsys@invoke\brl@pgfsys@invoke_matrix%
  \global\brl@pgf@invokemattoks{}%
}


\def\brl@pgf@matrix@endcell_plot{%
  \expandafter\xdef\csname brl@pgf@matrix_elem_\the\pgfmatrixcurrentrow _\the\pgfmatrixcurrentcolumn\endcsname{\the\brl@pgf@invokemattoks}%
}


\def\brl@pgf_matrix_plot{%
  \edef\brl@tmp{{\brl@pgf_beginscope%
      \brl@pgf_transformcm%
      {1.0}{0.0}%
      {0.0}{1.0}%
      {\the\@tempdimb}{\the\@tempdima}%
      \csname brl@pgf@matrix_elem_\the\@tempcnta _\the\@tempcntb\endcsname%
      \brl@pgf_endscope}}%
  \expandafter\pgfsysprotocol@literal\brl@tmp%
}%  




%%%%%%%%%%%%%%%%%%%%%%%%%%%%%%%%%%%%%%
%
% normalized normal vector
%
% input : \pgf@x, \pgf@y
% output : \brl@pgf@normal_x, \brl@pgf@normal_y
% \@tempdimb : sqrt[1+(b/a)^2]
% \@tempdimc : |a|
% \@tempdimb * \@tempdimc : sqrt[a^2+b^2]
%
\def\brl@pgf_normalize{% sqrt[1+x^2]
  \ifdim\pgf@x<\z@%
    \def\brl@pgf@xsign{-}%
  \else%
    \let\brl@pgf@xsign\empty%
  \fi%
  \ifdim\pgf@y<\z@%
    \def\brl@pgf@ysign{-}%
  \else%
    \let\brl@pgf@ysign\empty%
  \fi%
  \ifdim\brl@pgf@xsign\pgf@x>\brl@pgf@ysign\pgf@y% sqrt[a^2+b^2]=|a|sqrt[1+(b/a)^2]
    \@tempdimc\brl@pgf@xsign\pgf@x% |a|
    \@tempdima\dimexpr\pgf@y*\p@/\@tempdimc\relax% b/|a|
    \brl@pgf@xd\brl@pgf@xsign\brl@pgf@pointsize@d\p@% c sgn a
    \brl@pgf@yd\brl@pgf@pointsize@d\@tempdima% c b/|a|
    \brl@pgf_normalize_sqrt%
    \edef\brl@pgf@normal_x{\the\numexpr-\brl@pgf@yd/\@tempdimb\relax}%
    \edef\brl@pgf@normal_y{\the\numexpr\brl@pgf@xd/\@tempdimb\relax}%
  \else%
    %
    \ifdim\brl@pgf@ysign\pgf@y<0.01\p@% regard the vector as 0.
      \def\brl@pgf@normal_x{\z@}%
      \def\brl@pgf@normal_y{\z@}%
      \@tempdimb\p@%
      \@tempdimc\p@%
    \else%
      \@tempdimc\brl@pgf@ysign\pgf@y%
      \@tempdima\dimexpr\pgf@x*\p@/\@tempdimc\relax%
      \brl@pgf@xd\brl@pgf@pointsize@d\@tempdima%
      \brl@pgf@yd\brl@pgf@ysign\brl@pgf@pointsize@d\p@%
      \brl@pgf_normalize_sqrt%
      \edef\brl@pgf@normal_x{\the\numexpr-\brl@pgf@yd/\@tempdimb\relax}%
      \edef\brl@pgf@normal_y{\the\numexpr\brl@pgf@xd/\@tempdimb\relax}%
    \fi%
  \fi%
}


\def\brl@pgf_normalize_novec{% sqrt[1+x^2]
  \ifdim\pgf@x<\z@%
    \def\brl@pgf@xsign{-}%
  \else%
    \let\brl@pgf@xsign\empty%
  \fi%
  \ifdim\pgf@y<\z@%
    \def\brl@pgf@ysign{-}%
  \else%
    \let\brl@pgf@ysign\empty%
  \fi%
  \ifdim\brl@pgf@xsign\pgf@x>\brl@pgf@ysign\pgf@y% sqrt[a^2+b^2]=|a|sqrt[1+(b/a)^2]
    \@tempdimc\brl@pgf@xsign\pgf@x%
    \@tempdima\dimexpr\pgf@y*\p@/\pgf@x\relax%
    \brl@pgf_normalize_sqrt%
    \@tempdimb\expandafter\brl@@strip@pt\the\@tempdimc\@tempdimb%
  \else%
    \ifdim\brl@pgf@ysign\pgf@y<0.01\p@% regard the vector as 0.
      \@tempdimb\p@%
    \else%
      \@tempdimc\brl@pgf@ysign\pgf@y%
      \@tempdima\dimexpr\pgf@x*\p@/\pgf@y\relax%
      \brl@pgf_normalize_sqrt%
      \@tempdimb\expandafter\brl@@strip@pt\the\@tempdimc\@tempdimb%
    \fi%
  \fi%
}


%%%%%%%%%%%%%%%%%%%%%%%%%%%%%%%%%%%%%%%%%%%%%%%%%
%
% sqrt[1+x^2] for 0 <= x <= 1
%
% f[x_, a_] := (x + a / x) / 2;
% determine c so that the error of
% |f[1 + c x/2, 1 + x] / sqrt[1 + x] - 1| <= M for 0 <= x <= 1
% is minimized.
% M = 0.000052516
%
% numerical analysis implies c = 0.8575631489125308
% 1/c = 1.1660948832375695
% (c-2)/c = -1.332189766475139
% f[1 + c x/2, 1 + x] = (1/c) + 1/4 (2 + c x) + (c - 2)/(c (2 + c x)) 
%

\def\brl@pgf_normalize_sqrt{% sqrt[1 + x^2]
  \@tempdima\expandafter\brl@@strip@pt\the\@tempdima\@tempdima% x^2
  \@tempdima0.857563\@tempdima%
  \advance\@tempdima\tw@\p@% 2 + c x^2
  \@tempdimb\dimexpr1.16609\p@+0.25\@tempdima-1.33219\p@*\p@/\@tempdima\relax%
}

%%%%%%%%%%%%%%%%%%%%%%%%%%%%%%%%%%%%%%%%%%%%%%%%%
%
% 1/sqrt[1+x^2] for 0 <= x <= 1
%
% f[x_, a_] := (x + 1 /(a x)) / 2;
% determine c so that the error of
% |f[1 - c x/2, 1 + x] sqrt[1 + x] - 1| <= M for 0 <= x <= 1
% is minimized.
% M = 0.0005495844099567826
%
% numerical analysis implies c = 0.6319020829018348
% f[1 + c x/2, 1 + x] = 1/4 (2 - c x) + 1/(1 + x) (2 - c x)) 
%

% \def\brl@pgf_normalize_invsqrt{% sqrt[1 + x^2]
%   \@tempdima\expandafter\brl@@strip@pt\the\@tempdima\@tempdima% x^2
%   \@tempdimb\@tempdima%
%   \@tempdima-0.631902\@tempdima%
%   \advance\@tempdima\tw@\p@% 2 - c x^2
%   \advance\@tempdimb\p@% 1 + c x^2
%   \@tempdimb\dimexpr0.25\@tempdima\p@*\p@/\@tempdima\relax%
% }





% f[x_, a_] := (x + a / x) / 2;
% determine c so that the error of
% |f[1 + c x/2, 1 + x] - sqrt[1 + x]|
% is minimized.
%
% numerical analysis implies c = 0.8555572815447349
% 1/c = 1.1688288108476728
% (c-2)/c = -1.3376576216953455
% f[1 + c x/2, 1 + x] = (1/c) + 1/4 (2 + c x) + (c - 2)/(c (2 + c x)) 
%




\ifbrl@option@pindisplay%
  % Time-stamp: <2024-08-18 20:00:10 komori>
% Copyright (c) 2021 Yasushi Komori
% All rights reserved.

%%%%%%%%%%%%%%%%%%%%%%%%%%%%%%%%%%%%%%%%%%%%%%%%%%%%%%%%%%%%%%%%
%
% pgf for pindisplay
%

%%%%%%%%%%%%%%% 
% input:
%  \brl@pgf@currx, \brl@pgf@curry
%
% output:
%  determine whether this point is saved or not,
%  and save if needed.

\def\brl@pgf@oneletter{4}% total height
\def\brl@pgf_filldef@plotunit{1}%

\let\brl@pgf@leastdist\z@



%
% output:
% \brl@pgf@ptx, \brl@pgf@pty
% \box : sumiji letters for check
\protected\def\brl@pgf_hbox#1#2#3{%
  \brl@pgf@ptx#1\relax%
  \brl@pgf@pty#2\relax%
  %
  \count@\brl@pgf@minx\relax%
  \advance\count@\brl@pgf@ptx\relax%
  \ifodd\count@% a letter begins at odd.
    \advance\brl@pgf@ptx\@ne%
  \fi%
  %
  \setbox\pgf@hbox\hbox{\raise\numexpr\brl@pgf@pty*8\relax bp\hbox{\kern\numexpr\brl@pgf@ptx*8\relax bp #3}}%
  %
  \wd\pgf@hbox\z@%
  \ht\pgf@hbox\z@%
  \dp\pgf@hbox\z@%
  %
  \box\pgf@hbox%
}

\def\brl@pgf_drawpoint_init{%
  \edef\brl@pgf@transx{\the\brl@pgf@currx}%
  \edef\brl@pgf@transy{\the\brl@pgf@curry}%
}


\brl@pgf_filldef_save{4}{1}{1}{%
  {11}%
  {11}%
}
\brl@pgf_filldef_save{3}{1}{1}{%
  {01}%
  {11}%
}
\brl@pgf_filldef_save{2}{1}{1}{%
  {01}%
  {10}%
}
\brl@pgf_filldef_save{1}{1}{1}{%
  {00}%
  {10}%
}
\brl@pgf_filldef_save{0}{10000}{10000}{%
  {0}%
}


% grayscale

% 0 1 2 3 4

\def\brl@pgf_model_gray_#1\brl@end{%
  \count@\numexpr4-\dimexpr#1pt*4\relax/\p@\relax%
  % \showthe\count@
}

% gray := 0.3 · red + 0.59 · green + 0.11 · blue

\def\brl@pgf_model_rgb_#1,#2,#3\brl@end{%
  \count@\numexpr4-\dimexpr#1pt*3+#2pt*59+#3pt*11\relax/\p@/25\relax%
}

% gray := 1 − min{1, 0.3 · cyan + 0.59 · magenta + 0.11 · yellow + black}

\def\brl@pgf_model_cmyk_#1,#2,#3,#4\brl@end{%
  \count@\numexpr\dimexpr#1pt*3+#2pt*59+#3pt*11+#4pt\relax/\p@/25\relax%
  \ifnum4<\count@%
    \count@4\relax
  \fi%
}

% gray := 1 − (0.3 · cyan + 0.59 · magenta + 0.11 · yellow)

\def\brl@pgf_model_cmy_#1,#2,#3\brl@end{%
  \count@\numexpr\dimexpr#1pt*3+#2pt*59+#3pt*11\relax/\p@/25\relax%
}

\endinput


%%% Local Variables:
%%% mode: japanese-latex
%%% backup-each-save-backup : t
%%% TeX-master: t
%%% End:

\else%
  % Time-stamp: <2026-09-06 19:11:38 komori>
% Copyright (c) 2021 Yasushi Komor
% All rights reserved.

%%
%% for draw
%%

\def\pgfsys@setlinewidth#1{%
  \pgf@sys@bp{\brl@figure@linewidthscale\dimexpr#1\relax}% this line is patched.
  \pgfsysprotocol@literal{w}%
}

\newtoks\brl@pgf_filldef@toks% used only in draw


%% make boxs for ordinary texts and braille texts.
%% output : hbox \brl@pgf@nodeboxa, \brl@pgf@nodeboxb
%%


\expandafter\newbox\csname brl@pgf@nodetext_\string\pgfnodeparttextbox a\endcsname
\expandafter\newbox\csname brl@pgf@nodetext_\string\pgfnodeparttextbox b\endcsname

\def\pgfsys@hboxsynced#1{%
  \ifbrl@document@emptybox%
  \else%
    \pgfsys@begin@idscope%
    \pgfsys@beginscope%
    \pgf@sys@bp@correct\pgf@pt@x%
    \pgf@sys@bp@correct\pgf@pt@y%
    % 
    % 
    %
    \edef\brl@tmp{\pgf@pt@aa\brl@@space \pgf@pt@ab\brl@@space \pgf@pt@ba\brl@@space \pgf@pt@bb\brl@@space \pgf@sys@tonumber{\pgf@pt@x} \pgf@sys@tonumber{\pgf@pt@y} cm}%
    % 
    %
    \setbox\z@\hbox{%
      \box\csname brl@pgf@nodetext_\string#1a\endcsname%
    }%
    \wd\z@\z@%
    \ht\z@\z@%
    \dp\z@\z@%
    % 
    \setbox\z@\vbox{%
      \brl@@pdf_bcontent%
      \brl@@pdf_invoke{\brl@tmp}%
      \box\z@%
      \brl@@pdf_econtent%
    }%
    \wd\z@\z@%
    \ht\z@\z@%
    \dp\z@\z@%
    % 
    \setbox\@ne\hbox{%
      \box\csname brl@pgf@nodetext_\string#1b\endcsname%
    }%
    \wd\@ne\z@%
    \ht\@ne\z@%
    \dp\@ne\z@%
    % 
    \setbox\@ne\vbox{%
      \brl@@pdf_bcontent%
      \brl@@pdf_invoke{\brl@tmp}%
      \box\@ne%
      \brl@@pdf_econtent%
    }%
    \wd\@ne\z@%
    \ht\@ne\z@%
    \dp\@ne\z@%
    % 
    % 
    \brl@pgfsys@hboxsynced_sub%
    % 
    \pgfsys@endscope%
    \pgfsys@end@idscope%
  \fi%
}

\def\brl@pgfsys@hboxsynced_sub{%
  \global\setbox\brl@pgf@nodeboxa\hbox{\unhbox\brl@pgf@nodeboxa\box\z@}%
  \global\setbox\brl@pgf@nodeboxb\hbox{\unhbox\brl@pgf@nodeboxb\box\@ne}%
}


% \def\pgfsys@color@rgb@stroke#1#2#3{\pgfsysprotocol@literal{#1 #2 #3 RG}}
\def\pgfsys@color@rgb@fill#1#2#3{%
  \brl@pgf_color_setrgb{#1}{#2}{#3}%
  \edef\brl@tmp{{%
      /Pattern cs
      /PT\brl@pgf_filldef@pattern_name\brl@@space scn}}%
  \expandafter\pgfsysprotocol@literal\brl@tmp%
}
% \def\pgfsys@color@cmyk@stroke#1#2#3#4{\pgfsysprotocol@literal{#1 #2 #3 #4 K}}
\def\pgfsys@color@cmyk@fill#1#2#3#4{%
  \brl@pgf_color_setcmyk{#1}{#2}{#3}{#4}%
  \edef\brl@tmp{{%
      /Pattern cs
      /PT\brl@pgf_filldef@pattern_name\brl@@space scn}}%
  \expandafter\pgfsysprotocol@literal\brl@tmp%
}
% \def\pgfsys@color@cmy@stroke#1#2#3{\pgfsysprotocol@literal{#1 #2 #3 0 K}}
\def\pgfsys@color@cmy@fill#1#2#3{%
  \brl@pgf_color_setcmy{#1}{#2}{#3}%
  \edef\brl@tmp{{%
      /Pattern cs
      /PT\brl@pgf_filldef@pattern_name\brl@@space scn}}%
  \expandafter\pgfsysprotocol@literal\brl@tmp%
}
  % \def\pgfsys@color@gray@stroke#1{\pgfsysprotocol@literal{#1 G}}
\def\pgfsys@color@gray@fill#1{%
  \brl@pgf_color_setgray{#1}%
  \edef\brl@tmp{{%
      /Pattern cs
      /PT\brl@pgf_filldef@pattern_name\brl@@space scn}}%
  \expandafter\pgfsysprotocol@literal\brl@tmp%
}


\def\brl@pgf_color_setrgb#1#2#3{%
  \ifcsname brl@pgf_fill@pattern_def_rgb_#1,#2,#3\endcsname%
    \def\brl@pgf_filldef@pattern_name{rgb,#1,#2,#3}%
  \else%
    \ifcsname brl@pgf_filldef_save@rgb_#1,#2,#3\endcsname%
      \csname brl@pgf_filldef_save@rgb_#1,#2,#3\endcsname% exec user definition
    \else%
      \brl@pgf_model_rgb_#1,#2,#3\brl@end%
      %
      \ifcsname brl@pgf_fill@pattern_def_\the\count@\endcsname%
        \edef\brl@pgf_filldef@pattern_name{\the\count@}%
      \else%
        \csname brl@pgf_filldef_save@\the\count@\endcsname% exec predefinition
      \fi%
    \fi%
  \fi%
}%

\def\brl@pgf_color_setcmyk#1#2#3#4{%
  \ifcsname brl@pgf_fill@pattern_def_cmyk_#1,#2,#3,#4\endcsname%
    \def\brl@pgf_filldef@pattern_name{cmyk_#1,#2,#3,#4}%
  \else%
    \ifcsname brl@pgf_filldef_save@cmyk_#1,#2,#3,#4\endcsname%
      \csname brl@pgf_filldef_save@cmyk_#1,#2,#3,#4\endcsname% exec user definition
    \else%
      \brl@pgf_model_cmyk_#1,#2,#3,#4\brl@end%
      %
      \ifcsname brl@pgf_fill@pattern_def_\the\count@\endcsname%
        \edef\brl@pgf_filldef@pattern_name{\the\count@}%
      \else%
        \csname brl@pgf_filldef_save@\the\count@\endcsname% exec predefinition
      \fi%
    \fi%
  \fi%
}%

\def\brl@pgf_color_setcmy#1#2#3{%
  \ifcsname brl@pgf_fill@pattern_def_cmy_#1,#2,#3\endcsname%
    \def\brl@pgf_filldef@pattern_name{cmy_#1,#2,#3}%
  \else%
    \ifcsname brl@pgf_filldef_save@cmy_#1,#2,#3\endcsname%
      \csname brl@pgf_filldef_save@cmy_#1,#2,#3\endcsname% exec user definition
    \else%
      \brl@pgf_model_cmy_#1,#2,#3\brl@end%
      %
      \ifcsname brl@pgf_fill@pattern_def_\the\count@\endcsname%
        \edef\brl@pgf_filldef@pattern_name{\the\count@}%
      \else%
        \csname brl@pgf_filldef_save@\the\count@\endcsname% exec predefinition
      \fi%
    \fi%
  \fi%
}

\def\brl@pgf_color_setgray#1{%
  \ifcsname brl@pgf_fill@pattern_def_gray_#1\endcsname%
    \def\brl@pgf_filldef@pattern_name{gray_#1}%
  \else%
    \ifcsname brl@pgf_filldef_save@gray_#1\endcsname%
      \csname brl@pgf_filldef_save@gray_#1\endcsname% exec user definition
    \else%
      \brl@pgf_model_gray_#1\brl@end%
      %
      \ifcsname brl@pgf_fill@pattern_def_\the\count@\endcsname%
        \edef\brl@pgf_filldef@pattern_name{\the\count@}%
      \else%
        \csname brl@pgf_filldef_save@\the\count@\endcsname% exec predefinition
      \fi%
    \fi%
  \fi%
}

\ifbrl@lualatex
  % #1 #2 #3 #4 : r g b rg or x g X G (raw pdf code)
\def\brl@pgf_color_to_fill@g{g}%
\def\brl@pgf_color_to_fill@rg{rg}%
\def\brl@pgf_color_to_fill#1 #2 #3 #4\brl@end{%
  \def\brl@tmp{#2}%
  \ifx\brl@tmp\brl@pgf_color_to_fill@g% grayscale
    \brl@pgf_color_setgray{#1}%
  \else% rgb or cmyk
    \brl@pgf_color_to_fill_rgb#4\brl@end%
    \ifx\brl@tmp\brl@pgf_color_to_fill@rg% rgb
      \brl@pgf_color_setrgb{#1}{#2}{#3}%
    \else%
      \brl@pgf_color_setcmyk{#1}{#2}{#3}{\brl@tmp}%
    \fi%
  \fi%
}
\def\brl@pgf_color_to_fill_rgb#1 #2\brl@end{%
  \def\brl@tmp{#1}%
}

\else%
  
\def\brl@pgf_color_to_fill#1 {%
  \csname brl@pgf_color_to_fill_#1\endcsname%
}
\def\brl@pgf_color_to_fill_rgb#1 #2 #3\brl@end{%
  \brl@pgf_color_setrgb{#1}{#2}{#3}%
}

\def\brl@pgf_color_to_fill_cmyk#1 #2 #3 #4\brl@end{%
  \brl@pgf_color_setcmyk{#1}{#2}{#3}{#4}%
}

\def\brl@pgf_color_to_fill_cmy#1 #2 #3\brl@end{%
  \brl@pgf_color_setcmy{#1}{#2}{#3}%
}

\def\brl@pgf_color_to_fill_gray#1\brl@end{%
  \brl@pgf_color_setgray{#1}%
}

\fi%


\def\brl@pgfsys@set@color{%
  \expandafter\brl@pgf_color_to_fill\current@color\brl@end%
  \edef\brl@tmp{{%
      q
      /Pattern cs
      /PT\brl@pgf_filldef@pattern_name\brl@@space scn}}%
  \expandafter\brl@@pdf_invoke\brl@tmp%
  \aftergroup\brl@pgfsys@reset@color%
}
\def\brl@pgfsys@reset@color{\brl@@pdf_invoke{Q}}

%%%%%%%%%%%%%%%%%%%%%%%%%%%%%%%%%%%%%%%%%%%%%%%%%%%%%%%%%%%%%%%

\def\brl@pgf_draw_filldef#1#2#3#4{%
  \edef\brl@pgf_filldef@pattern_name{#1}%
  \global\expandafter\let\csname brl@pgf_fill@pattern_def_\brl@pgf_filldef@pattern_name\endcsname\empty% switch, \@undefined or \empty
  \brl@pgf_filldef@curry\z@%
  \def\brl@pgf_filldef@maxx{-1}%
  \edef\brl@pgf_filldef@minstepx{\the\dimexpr\brl@pgf_filldef@drawunit pt*#2\relax}%
  \edef\brl@pgf_filldef@minstepy{\the\dimexpr\brl@pgf_filldef@drawunit pt*#2\relax}%
  % 
  \brl@pgf_draw_filldef_count#4\brl@end% to reverse the order, count the number of patterns
  \edef\brl@pgf_filldef@maxy{\the\brl@pgf_filldef@curry}%
  %
  \ifnum#2=10000\relax%
    \def\brl@pgf_filldef@drawmaxx{100}%
    \def\brl@pgf_filldef@drawmaxy{100}%
    %
  \else
    %
    \edef\brl@pgf_filldef@drawmaxx{\expandafter\brl@@strip@pt\the\dimexpr\brl@pgf_filldef@minstepx*\brl@pgf_filldef@maxx\relax}%
    \edef\brl@pgf_filldef@drawmaxy{\expandafter\brl@@strip@pt\the\dimexpr\brl@pgf_filldef@minstepy*\brl@pgf_filldef@maxy\relax}%
  \fi%
  %
  \let\brl@pgf_filldef@pattern\empty%
  \brl@pgf_draw_filldef_pattern#4%
}

\def\brl@pgf_draw_filldef_count#1{%
  \def\brl@tmp{#1}%
  \ifx\brl@tmp\brl@@end%
  \else%
    \brl@pgf_filldef@currx\z@%
    \brl@pgf_draw_filldef_count_x#1\brl@end%
    \advance\brl@pgf_filldef@curry\@ne%
    \expandafter\brl@pgf_draw_filldef_count%
  \fi%
}

\def\brl@pgf_draw_filldef_count_x#1{%
  \def\brl@tmp{#1}%
  \ifx\brl@tmp\brl@@end%
    \ifnum\brl@pgf_filldef@maxx<\brl@pgf_filldef@currx%
      \edef\brl@pgf_filldef@maxx{\the\brl@pgf_filldef@currx}%
    \fi%
  \else%
    \advance\brl@pgf_filldef@currx\@ne%
    \expandafter\brl@pgf_draw_filldef_count_x%
  \fi%
}



\def\brl@pgf_draw_filldef_pattern#1{%
  %
  \advance\brl@pgf_filldef@curry\m@ne%
  \brl@pgf_filldef@currx\z@%
  % 
  \brl@pgf_draw_filldef_pattern_loop#1\brl@end%
  \ifnum\brl@pgf_filldef@curry>\z@%
    \expandafter\brl@pgf_draw_filldef_pattern%
  \else%
    \expandafter\brl@pgf_draw_filldef_pattern_end%
  \fi%
}

\def\brl@pgf_draw_filldef_pattern_loop#1{%
  \def\brl@pgf_filldef@currsize{#1}%
  \ifx\brl@pgf_filldef@currsize\brl@@end%
    %
    % end
    %
  \else%
    %
    % loop
    %
    \ifnum\brl@pgf_filldef@currsize>\z@%
      %
      % \brl@pgf_filldef@currx
      % \brl@pgf_filldef@curry
      % \brl@pgf_filldef@currsize
      %
      \edef\brl@pgf_filldef@drawx{\expandafter\brl@@strip@pt\the\dimexpr\brl@pgf_filldef@minstepx*\brl@pgf_filldef@currx\relax}%
      \edef\brl@pgf_filldef@drawy{\expandafter\brl@@strip@pt\the\dimexpr\brl@pgf_filldef@minstepy*\brl@pgf_filldef@curry\relax}%
      % 
      \brl@pgf_draw_filldef_pattern_set%
    \fi%
    %
    % 
    \advance\brl@pgf_filldef@currx\@ne%
    %
    \expandafter\brl@pgf_draw_filldef_pattern_loop%
  \fi%
}

\def\brl@pgf_draw_filldef_pattern_set{%
  \edef\brl@pgf_filldef@pattern{\brl@pgf_filldef@pattern%
    q \brl@pgf_filldef@currsize\brl@@space 0 0 \brl@pgf_filldef@currsize\brl@@space%
    \brl@pgf_filldef@drawx\brl@@space%
    \brl@pgf_filldef@drawy\brl@@space cm\brl@@space%
    \brl@pgf_filldef@drawpattern\brl@@space%
    Q }%
  \ifdim\brl@pgf_filldef@drawx\p@=\z@%
    \edef\brl@pgf_filldef@pattern{\brl@pgf_filldef@pattern%
      q \brl@pgf_filldef@currsize\brl@@space 0 0 \brl@pgf_filldef@currsize\brl@@space%
      \brl@pgf_filldef@drawmaxx\brl@@space%
      \brl@pgf_filldef@drawy\brl@@space cm\brl@@space%
      \brl@pgf_filldef@drawpattern\brl@@space%
      Q }%
    \ifdim\brl@pgf_filldef@drawy\p@=\z@%
      \edef\brl@pgf_filldef@pattern{\brl@pgf_filldef@pattern%
        q \brl@pgf_filldef@currsize\brl@@space 0 0 \brl@pgf_filldef@currsize\brl@@space%
        \brl@pgf_filldef@drawmaxx\brl@@space%
        \brl@pgf_filldef@drawmaxy\brl@@space cm\brl@@space%
        \brl@pgf_filldef@drawpattern\brl@@space%
        Q q \brl@pgf_filldef@currsize\brl@@space 0 0 \brl@pgf_filldef@currsize\brl@@space%
        \brl@pgf_filldef@drawx\brl@@space%
        \brl@pgf_filldef@drawmaxy\brl@@space cm\brl@@space%
        \brl@pgf_filldef@drawpattern\brl@@space%
        Q }%
    \fi%
  \else%
    \ifdim\brl@pgf_filldef@drawy\p@=\z@%
      \edef\brl@pgf_filldef@pattern{\brl@pgf_filldef@pattern%
        q \brl@pgf_filldef@currsize\brl@@space 0 0 \brl@pgf_filldef@currsize\brl@@space%
        \brl@pgf_filldef@drawx\brl@@space%
        \brl@pgf_filldef@drawmaxy\brl@@space cm\brl@@space%
        \brl@pgf_filldef@drawpattern\brl@@space%
        Q }%
    \fi%
  \fi%
}%




\ifbrl@lualatex

\def\brl@pgf_draw_filldef_pattern_end{%
  \edef\brl@tmp{{%
      \the\brl@pgf_filldef@toks%
      \immediate\pdfextension obj stream attr{
        /Type /Pattern
        /PatternType 1
        /PaintType 1
        /TilingType 3
        /BBox [0 0 \brl@pgf_filldef@drawmaxx\brl@@space\brl@pgf_filldef@drawmaxy]
        /XStep \brl@pgf_filldef@drawmaxx\brl@@space
        /YStep \brl@pgf_filldef@drawmaxy\brl@@space
        /Resources << >>
        /Matrix [1 0 0 1 0 0]
      }{\brl@pgf_filldef@pattern}%
      \xdef\noexpand\brl@@patterns_list{\noexpand\brl@@patterns_list\brl@@space%
        /PT\brl@pgf_filldef@pattern_name\brl@@space\noexpand\pdffeedback lastobj 0 R}%
    }}%
  \global\brl@pgf_filldef@toks\brl@tmp%
}

\else%

\def\brl@pgf_draw_filldef_pattern_end{%
  \edef\brl@tmp{{%
      \the\brl@pgf_filldef@toks%
      \special{pdf:stream @fillpattern\brl@pgf_filldef@pattern_name\brl@@space
        (\brl@pgf_filldef@pattern)
        <<
        /Type /Pattern
        /PatternType 1
        /PaintType 1
        /TilingType 3
        /BBox [0 0 \brl@pgf_filldef@drawmaxx\brl@@space\brl@pgf_filldef@drawmaxy]
        /XStep \brl@pgf_filldef@drawmaxx\brl@@space
        /YStep \brl@pgf_filldef@drawmaxy\brl@@space
        /Resources << >>
        /Matrix [1 0 0 1 0 0]
        >>
      }%
    }}%
  \global\brl@pgf_filldef@toks\brl@tmp%
  \xdef\brl@@patterns_list{\brl@@patterns_list\brl@@space%
    /PT\brl@pgf_filldef@pattern_name\brl@@space @fillpattern\brl@pgf_filldef@pattern_name}% register here.
}

\fi%


%%%
%
% \matrix

\def\brl@pgf_draw_matrix_init{%
  \let\brl@pgfsys@hboxsynced_sub\brl@pgfsys@hboxsynced_matrix%
}

\def\brl@pgfsys@hboxsynced_matrix{%
  \edef\brl@tmp{brl@pgf@matrix_elem_\the\pgfmatrixcurrentrow _\the\pgfmatrixcurrentcolumn}%
  % create node for each element in a matrix.
  \global\setbox\csname\brl@tmp a\endcsname\hbox{\unhbox\csname\brl@tmp a\endcsname\box\z@}%
  \global\setbox\csname\brl@tmp b\endcsname\hbox{\unhbox\csname\brl@tmp b\endcsname\box\@ne}%
}

\def\brl@pgf@matrix@step@column_sub{%
  % prepare box
  %
  \edef\brl@tmp{brl@pgf@matrix_elem_\the\pgfmatrixcurrentrow _\the\pgfmatrixcurrentcolumn}%
  \ifcsname\brl@tmp p\endcsname%
  \else%
    \expandafter\newbox\csname\brl@tmp p\endcsname% picture
    \expandafter\newbox\csname\brl@tmp a\endcsname% nodeboxa
    \expandafter\newbox\csname\brl@tmp b\endcsname% nodeboxb
  \fi%
}

%
\def\brl@pgf@matrix@endcell_sub{%
  \setbox\z@\lastbox%
  \global\setbox\csname brl@pgf@matrix_elem_\the\pgfmatrixcurrentrow _\the\pgfmatrixcurrentcolumn p\endcsname\hbox{\copy\z@}%
  \box\z@%
}


\def\brl@pgf_matrix_sub{%
  \setbox\@ne\hbox{%
      \unhbox\csname brl@pgf@matrix_elem_\the\@tempcnta _\the\@tempcntb p\endcsname%
    }%
  \wd\@ne\z@%
  \ht\@ne\z@%
  \dp\@ne\z@%
  %
  \setbox\z@\hbox{\hskip\@tempdimb%
    \raise\@tempdima\hbox{%
      \unhcopy\@ne%
      \unhbox\csname brl@pgf@matrix_elem_\the\@tempcnta _\the\@tempcntb a\endcsname%
    }}%
  \wd\z@\z@%
  \ht\z@\z@%
  \dp\z@\z@%
  \global\setbox\brl@pgf@nodeboxa\hbox{\unhbox\brl@pgf@nodeboxa\box\z@}%
  %
  \setbox\z@\hbox{\hskip\@tempdimb%
    \raise\@tempdima\hbox{%
      \unhbox\@ne%
      \unhbox\csname brl@pgf@matrix_elem_\the\@tempcnta _\the\@tempcntb b\endcsname%
    }}%
  \wd\z@\z@%
  \ht\z@\z@%
  \dp\z@\z@%
  \global\setbox\brl@pgf@nodeboxb\hbox{\unhbox\brl@pgf@nodeboxb\box\z@}%
}%  



%%% Local Variables:
%%% mode: LaTeX
%%% TeX-master: t
%%% End:
  
  % Time-stamp: <2024-10-10 15:00:14 komori>
% Copyright (c) 2021 Yasushi Komori
% All rights reserved.

\def\brl@pgf@oneletter{31}% total height
\def\brl@pgf@dbcht{14}% dbc height
%% 14 is the same as the highest point of a letter
%% \@namedef{brl@bc_plot@x0}{0}
%% \@namedef{brl@bc_plot@y0}{14}

\def\brl@pgf@leastdist{\csname brl@pgf@linestep\brl@pgf@pointsize\endcsname\relax}



%%
%%


%
% output:
% \brl@pgf@ptx, \brl@pgf@pty
% \box : sumiji letters for check
\protected\def\brl@pgf_hbox#1#2#3{%
  \brl@pgf@ptx#1\relax%
  \brl@pgf@pty#2\relax%
  %
  \setbox\pgf@hbox\hbox{\raise\brl@pgf@pty bp\hbox{\brl@orig@hskip\brl@pgf@ptx bp #3}}%
  %
  \wd\pgf@hbox\z@%
  \ht\pgf@hbox\z@%
  \dp\pgf@hbox\z@%
  %
  \box\pgf@hbox%
}

\def\brl@pgf_drawpoint_init{%
  \brl@pgf_orient%
  \ifnum\brl@pgf@pointsize=\z@% if 0, then nbd is not removed.
  \else%
  % the radius of the removal
    \edef\brl@pgf_drawpoint@maxy{\the\numexpr\brl@pgf@pointsize+\brl@pgf@radius\relax}%
    \brl@pgf@dy-\brl@pgf_drawpoint@maxy\relax%
    \expandafter\brl@pgf_drawpoint_remove_nbd% remove all the neighboring points
  \fi%
}


\def\brl@pgf_drawpoint_remove_nbd{%
  \edef\brl@pgf_drawpoint@curry{_\the\numexpr\brl@pgf@dy+\brl@pgf@transy\relax}%
  \ifcsname brl@pgf@minxaty\brl@pgf_drawpoint@curry\endcsname%
    \csname brl@pgf_drawpoint@x_\brl@pgf_drawpoint@maxy _\the\brl@pgf@dy\endcsname%
  \fi%
  \ifnum\brl@pgf_drawpoint@maxy>\brl@pgf@dy%
    \advance\brl@pgf@dy\@ne%
    \expandafter\brl@pgf_drawpoint_remove_nbd%
  \fi%
}

\def\brl@pgf_orient_portrait{%
  \edef\brl@pgf@transx{\the\brl@pgf@currx}%
  \edef\brl@pgf@transy{\the\brl@pgf@curry}%
}

\def\brl@pgf_orient_landscape{%
  \edef\brl@pgf@transx{\the\numexpr-\brl@pgf@curry\relax}%
  \edef\brl@pgf@transy{\the\brl@pgf@currx}%
}


% input: \pgf@x, \pgf@y, \brl@pgf@prevx, \brl@pgf@prevy
%
% determine whether the point is set or not.
%
%%
%% shape of the neighbor
%%
%
% #1 : maxy
% #2 : y
% #3 : width

\def\brl@tmp{{%
    \global\expandafter\let\csname brl@pgf@size_\the\brl@pgf@dx\brl@pgf_drawpoint@curry\endcsname\@undefined%
  }}%

\def\brl@pgf_drawpointdef#1#2#3{%
  \def\brl@pgf_drawpointdef@width{#3}%
  \count@-#3\relax%
  \brl@def@toks{}%
  \brl@pgf_drawpointdef_loop%
  \expandafter\edef\csname brl@pgf_drawpoint@x_#1_#2\endcsname{\unexpanded{\brl@pgf@dx\numexpr\brl@pgf@transx-#3\relax}\the\brl@def@toks}%
}

\def\brl@pgf_drawpointdef_loop{%
  \expandafter\brl@def_appendtoks\brl@tmp%
  % 
  \ifnum\brl@pgf_drawpointdef@width>\count@%
    \brl@def_appendtoks{\advance\brl@pgf@dx\@ne}%
    \advance\count@\@ne%
    \expandafter\brl@pgf_drawpointdef_loop%
  \fi%
}


% >38
%　　　　　■■ 72 61 52 45 40 37 36
%　　　■■■■ 61 50 41 34 29 26 25
%　　■■■■■ 52 41 32 25 20 17 16
%　■■■■■■ 45 34 25 18 13 10  9
%　■■■■■■ 40 29 20 13  8  5  4
%■■■■■■■ 37 26 17 10  5  2  1
%■■■■■■■ 36 25 16  9  4  1  0


\brl@pgf_drawpointdef{6}{6}{1}
\brl@pgf_drawpointdef{6}{5}{3}
\brl@pgf_drawpointdef{6}{4}{4}
\brl@pgf_drawpointdef{6}{3}{5}
\brl@pgf_drawpointdef{6}{2}{5}
\brl@pgf_drawpointdef{6}{1}{6}
\brl@pgf_drawpointdef{6}{0}{6}
\brl@pgf_drawpointdef{6}{-1}{6}
\brl@pgf_drawpointdef{6}{-2}{5}
\brl@pgf_drawpointdef{6}{-3}{5}
\brl@pgf_drawpointdef{6}{-4}{4}
\brl@pgf_drawpointdef{6}{-5}{3}
\brl@pgf_drawpointdef{6}{-6}{1}

% >26
%　　　　■■ 50 41 34 29 26 25
%　　■■■■ 41 32 25 20 17 16
%　■■■■■ 34 25 18 13 10  9
%　■■■■■ 29 20 13  8  5  4
%■■■■■■ 26 17 10  5  2  1
%■■■■■■ 25 16  9  4  1  0

\brl@pgf_drawpointdef{5}{5}{1}
\brl@pgf_drawpointdef{5}{4}{3}
\brl@pgf_drawpointdef{5}{3}{4}
\brl@pgf_drawpointdef{5}{2}{4}
\brl@pgf_drawpointdef{5}{1}{5}
\brl@pgf_drawpointdef{5}{0}{5}
\brl@pgf_drawpointdef{5}{-1}{5}
\brl@pgf_drawpointdef{5}{-2}{4}
\brl@pgf_drawpointdef{5}{-3}{4}
\brl@pgf_drawpointdef{5}{-4}{3}
\brl@pgf_drawpointdef{5}{-5}{1}

% >18
%　　　■■ 32 25 20 17 16
%　■■■■ 25 18 13 10  9
%　■■■■ 20 13  8  5  4
%■■■■■ 17 10  5  2  1
%■■■■■ 16  9  4  1  0

\brl@pgf_drawpointdef{4}{4}{1}
\brl@pgf_drawpointdef{4}{3}{3}
\brl@pgf_drawpointdef{4}{2}{3}
\brl@pgf_drawpointdef{4}{1}{4}
\brl@pgf_drawpointdef{4}{0}{4}
\brl@pgf_drawpointdef{4}{-1}{4}
\brl@pgf_drawpointdef{4}{-2}{3}
\brl@pgf_drawpointdef{4}{-3}{3}
\brl@pgf_drawpointdef{4}{-4}{1}

% >10
%　　■■ 18 13 10  9
%　■■■ 13  8  5  4
%■■■■ 10  5  2  1
%■■■■  9  4  1  0

\brl@pgf_drawpointdef{3}{3}{1}
\brl@pgf_drawpointdef{3}{2}{2}
\brl@pgf_drawpointdef{3}{1}{3}
\brl@pgf_drawpointdef{3}{0}{3}
\brl@pgf_drawpointdef{3}{-1}{3}
\brl@pgf_drawpointdef{3}{-2}{2}
\brl@pgf_drawpointdef{3}{-3}{1}

% >5
%　■■ 8  5  4
%■■■ 5  2  1
%■■■ 4  1  0

\brl@pgf_drawpointdef{2}{2}{1}
\brl@pgf_drawpointdef{2}{1}{2}
\brl@pgf_drawpointdef{2}{0}{2}
\brl@pgf_drawpointdef{2}{-1}{2}
\brl@pgf_drawpointdef{2}{-2}{1}


%%
%% line shape
%%
%% height and depth = radius+size
%%
\def\brl@pgf@step_sparse{%
  \def\brl@pgf@radius{2}%
  \@namedef{brl@pgf@linestep1}{25}%
  \@namedef{brl@pgf@linestep2}{36}%
  \@namedef{brl@pgf@linestep3}{49}%
  \@namedef{brl@pgf@linestep4}{25}%
  \@namedef{brl@pgf@linestep5}{36}%
  \@namedef{brl@pgf@linestep6}{49}%
  \@namedef{brl@pgf@linestep7}{25}%
  \@namedef{brl@pgf@linestep8}{36}%
  \@namedef{brl@pgf@linestep9}{49}%
  \def\brl@pgf_filldef@plotunit{7}%
}

\def\brl@pgf@step_loose{%
  \def\brl@pgf@radius{2}%
  \@namedef{brl@pgf@linestep1}{16}%
  \@namedef{brl@pgf@linestep2}{25}%
  \@namedef{brl@pgf@linestep3}{36}%
  \@namedef{brl@pgf@linestep4}{16}%
  \@namedef{brl@pgf@linestep5}{25}%
  \@namedef{brl@pgf@linestep6}{36}%
  \@namedef{brl@pgf@linestep7}{16}%
  \@namedef{brl@pgf@linestep8}{25}%
  \@namedef{brl@pgf@linestep9}{36}%
  \def\brl@pgf_filldef@plotunit{6}%
}

\def\brl@pgf@step_natural{%
  \def\brl@pgf@radius{2}%
  \@namedef{brl@pgf@linestep1}{10}%
  \@namedef{brl@pgf@linestep2}{18}%
  \@namedef{brl@pgf@linestep3}{26}%
  \@namedef{brl@pgf@linestep4}{10}%
  \@namedef{brl@pgf@linestep5}{18}%
  \@namedef{brl@pgf@linestep6}{26}%
  \@namedef{brl@pgf@linestep7}{10}%
  \@namedef{brl@pgf@linestep8}{18}%
  \@namedef{brl@pgf@linestep9}{26}%
  \def\brl@pgf_filldef@plotunit{6}%
}
% 「多めに探す」でいくつか
\def\brl@pgf@step_dense{%
  \def\brl@pgf@radius{1}%
  \@namedef{brl@pgf@linestep1}{9}%
  \@namedef{brl@pgf@linestep2}{16}%
  \@namedef{brl@pgf@linestep3}{25}%
  \@namedef{brl@pgf@linestep4}{9}%
  \@namedef{brl@pgf@linestep5}{16}%
  \@namedef{brl@pgf@linestep6}{25}%
  \@namedef{brl@pgf@linestep7}{9}%
  \@namedef{brl@pgf@linestep8}{16}%
  \@namedef{brl@pgf@linestep9}{25}%
  \def\brl@pgf_filldef@plotunit{5}%
}
% 「標準レベル」で 0, 「多めに探す」で 全部 (ぎりぎり)
\def\brl@pgf@step_continuous{%
  \def\brl@pgf@radius{1}%
  \@namedef{brl@pgf@linestep1}{5}% 
  \@namedef{brl@pgf@linestep2}{10}%
  \@namedef{brl@pgf@linestep3}{18}%
  \@namedef{brl@pgf@linestep4}{5}% 
  \@namedef{brl@pgf@linestep5}{10}%
  \@namedef{brl@pgf@linestep6}{18}%
  \@namedef{brl@pgf@linestep7}{5}% 
  \@namedef{brl@pgf@linestep8}{10}%
  \@namedef{brl@pgf@linestep9}{18}%
  \def\brl@pgf_filldef@plotunit{5}%
}






%%%%%%%%%%%%%%%%%%%%%%%%%%%%%%%%%%%%%%%%%%%%%%%%%
%%
%% grayscale
%%
%%%%%%%%%%%%%%%%%%%%%%%%%%%%%%%%%%%%%%%%%%%%%%%%%

\brl@pgf_filldef_save{50}{1}{1}{%
  {3333}%
  {3333}%
  {3333}%
  {3333}%
}
\brl@pgf_filldef_save{49}{1}{1}{%
  {1333}%
  {3333}%
  {3333}%
  {3333}%
}
\brl@pgf_filldef_save{48}{1}{1}{%
  {1333}%
  {3333}%
  {3313}%
  {3333}%
}
\brl@pgf_filldef_save{47}{1}{1}{%
  {1313}%
  {3333}%
  {3313}%
  {3333}%
}
\brl@pgf_filldef_save{46}{1}{1}{%
  {1313}%
  {3333}%
  {1313}%
  {3333}%
}
\brl@pgf_filldef_save{45}{1}{1}{%
  {1313}%
  {3133}%
  {1313}%
  {3333}%
}
\brl@pgf_filldef_save{44}{1}{1}{%
  {1313}%
  {3133}%
  {1313}%
  {3331}%
}
\brl@pgf_filldef_save{43}{1}{1}{%
  {1313}%
  {3131}%
  {1313}%
  {3331}%
}
\brl@pgf_filldef_save{42}{1}{1}{%
  {1313}%
  {3131}%
  {1313}%
  {3131}%
}
\brl@pgf_filldef_save{41}{1}{1}{%
  {1113}%
  {3131}%
  {1313}%
  {3131}%
}
\brl@pgf_filldef_save{40}{1}{1}{%
  {1113}%
  {3131}%
  {1311}%
  {3131}%
}
\brl@pgf_filldef_save{39}{1}{1}{%
  {1111}%
  {3131}%
  {1311}%
  {3131}%
}
\brl@pgf_filldef_save{38}{1}{1}{%
  {1111}%
  {3131}%
  {1111}%
  {3131}%
}
\brl@pgf_filldef_save{37}{1}{1}{%
  {1111}%
  {1131}%
  {1111}%
  {3131}%
}
\brl@pgf_filldef_save{36}{1}{1}{%
  {1111}%
  {1131}%
  {1111}%
  {3111}%
}
\brl@pgf_filldef_save{35}{1}{1}{%
  {1111}%
  {1111}%
  {1111}%
  {3111}%
}
%%%%%%%%%%%%%%%%%%%%%%%%%%%%%%%%%%%%%%%%%%%%
\brl@pgf_filldef_save{34}{1}{1}{%
  {2222}%
  {2222}%
  {2222}%
  {2222}%
}
\brl@pgf_filldef_save{33}{1}{1}{%
  {0222}%
  {2222}%
  {2222}%
  {2222}%
}
\brl@pgf_filldef_save{32}{1}{1}{%
  {0222}%
  {2222}%
  {2202}%
  {2222}%
}
\brl@pgf_filldef_save{31}{1}{1}{%
  {0202}%
  {2222}%
  {2202}%
  {2222}%
}
\brl@pgf_filldef_save{30}{1}{1}{%
  {0202}%
  {2222}%
  {0202}%
  {2222}%
}
\brl@pgf_filldef_save{29}{1}{1}{%
  {0202}%
  {2022}%
  {0202}%
  {2222}%
}
\brl@pgf_filldef_save{28}{1}{1}{%
  {0202}%
  {2022}%
  {0202}%
  {2220}%
}
\brl@pgf_filldef_save{27}{1}{1}{%
  {0202}%
  {2020}%
  {0202}%
  {2220}%
}
\brl@pgf_filldef_save{26}{1}{1}{%
  {0202}%
  {2020}%
  {0202}%
  {2020}%
}
\brl@pgf_filldef_save{25}{1}{1}{%
  {0002}%
  {2020}%
  {0202}%
  {2020}%
}
\brl@pgf_filldef_save{24}{1}{1}{%
  {0002}%
  {2020}%
  {0200}%
  {2020}%
}
\brl@pgf_filldef_save{23}{1}{1}{%
  {0000}%
  {2020}%
  {0200}%
  {2020}%
}
\brl@pgf_filldef_save{22}{1}{1}{%
  {0000}%
  {2020}%
  {0000}%
  {2020}%
}
\brl@pgf_filldef_save{21}{1}{1}{%
  {0000}%
  {0020}%
  {0000}%
  {2020}%
}
\brl@pgf_filldef_save{20}{1}{1}{%
  {0000}%
  {0020}%
  {0000}%
  {2000}%
}
\brl@pgf_filldef_save{19}{1}{1}{%
  {0000}%
  {0000}%
  {0000}%
  {2000}%
}

%%%%%%%%%%%%%%%%%%%%%%%%%%%%%%%%%%%%%%%%%%%%
\brl@pgf_filldef_save{18}{1}{1}{%
  {1111}%
  {1111}%
  {1111}%
  {1111}%
}
\brl@pgf_filldef_save{17}{1}{1}{%
  {0111}%
  {1111}%
  {1111}%
  {1111}%
}
\brl@pgf_filldef_save{16}{1}{1}{%
  {0111}%
  {1111}%
  {1101}%
  {1111}%
}
\brl@pgf_filldef_save{15}{1}{1}{%
  {0101}%
  {1111}%
  {1101}%
  {1111}%
}
\brl@pgf_filldef_save{14}{1}{1}{%
  {0101}%
  {1111}%
  {0101}%
  {1111}%
}
\brl@pgf_filldef_save{13}{1}{1}{%
  {0101}%
  {1011}%
  {0101}%
  {1111}%
}
\brl@pgf_filldef_save{12}{1}{1}{%
  {0101}%
  {1011}%
  {0101}%
  {1110}%
}
\brl@pgf_filldef_save{11}{1}{1}{%
  {0101}%
  {1010}%
  {0101}%
  {1110}%
}
\brl@pgf_filldef_save{10}{1}{1}{%
  {0101}%
  {1010}%
  {0101}%
  {1010}%
}
\brl@pgf_filldef_save{9}{1}{1}{%
  {0001}%
  {1010}%
  {0101}%
  {1010}%
}
\brl@pgf_filldef_save{8}{1}{1}{%
  {0001}%
  {1010}%
  {0100}%
  {1010}%
}
\brl@pgf_filldef_save{7}{1}{1}{%
  {0000}%
  {1010}%
  {0100}%
  {1010}%
}
\brl@pgf_filldef_save{6}{1}{1}{%
  {0000}%
  {1010}%
  {0000}%
  {1010}%
}
\brl@pgf_filldef_save{5}{1}{1}{%
  {0000}%
  {0010}%
  {0000}%
  {1010}%
}
\brl@pgf_filldef_save{4}{1}{1}{%
  {0000}%
  {0010}%
  {0000}%
  {1000}%
}
\brl@pgf_filldef_save{3}{1}{1}{%
  {0000}%
  {0000}%
  {0000}%
  {1000}%
}
\brl@pgf_filldef_save{2}{1}{1}{%
  {0000}%
  {0000}%
  {0000}%
  {1000}%
}
\brl@pgf_filldef_save{1}{1}{1}{%
  {0000}%
  {0000}%
  {0000}%
  {1000}%
}
\brl@pgf_filldef_save{0}{10000}{10000}{%
  {0}%
}




%
% texdoc xcolor
%
% grayscale

\def\brl@pgf_model_gray_#1\brl@end{%
  \count@\numexpr50-\dimexpr#1pt*50\relax/\p@\relax%
}

% gray := 0.3 · red + 0.59 · green + 0.11 · blue

\def\brl@pgf_model_rgb_#1,#2,#3\brl@end{%
  \count@\numexpr50-\dimexpr#1pt*3+#2pt*59+#3pt*11\relax/\p@/2\relax%
}

% gray := 1 − min{1, 0.3 · cyan + 0.59 · magenta + 0.11 · yellow + black}

\def\brl@pgf_model_cmyk_#1,#2,#3,#4\brl@end{%
  \count@\numexpr\dimexpr#1pt*3+#2pt*59+#3pt*11+#4pt\relax/\p@/2\relax%
  \ifnum50<\count@%
    \count@50\relax%
  \fi%
}

% gray := 1 − (0.3 · cyan + 0.59 · magenta + 0.11 · yellow)

\def\brl@pgf_model_cmy_#1,#2,#3\brl@end{%
  \count@\numexpr\dimexpr#1pt*3+#2pt*59+#3pt*11\relax/\p@/2\relax%
}


\endinput


%%% Local Variables:
%%% mode: japanese-latex
%%% backup-each-save-backup : t
%%% TeX-master: t
%%% End:

\fi%

\endinput


%%% Local Variables:
%%% mode: japanese-latex
%%% backup-each-save-backup : t
%%% TeX-master: t
%%% End:
